\unit{Exposé 221 -- Les schémas de Hilbert}
\sid{221}
\src{249}
\seminar{13e année, 1960/61, n\textsuperscript{o} 221\\Mai 1961}
\paperhead{TECHNIQUES DE CONSTRUCTION ET THÉORÈMES D'EXISTENCE EN GÉOMÉTRIE ALGÉBRIQUE\\
IV : LES SCHÉMAS DE HILBERT}%
  {par Alexander GROTHENDIECK}

\phantomsection\addcontentsline{toc}{subsection}{Introduction}
\noindent INTRODUCTION. - Les techniques exposées dans [2], I et II, étaient pour l'essentiel
indépendantes de toute hypothèse projective sur les schémas envisagés.
Malheureusement, elles ne permettent pas à l'heure actuelle de résoudre les
problèmes d'existence posés dans l'exposé II de [2]. Dans le présent exposé et
le suivant, nous allons résoudre ces problèmes moyennant des hypothèses
projectives. Les techniques employées sont typiquement projectives, et ne font
pratiquement pas usage des résultats de I et II de [2]. Ici, nous allons
construire les "schémas de Hilbert" destinés à remplacer l'utilisation des
coordonnées de Chow, comme il a été dit dans [2], II, n\textsuperscript{o} 2.
Dans l'exposé suivant, la théorie de passage au quotient dans les schémas
développée dans III, jointe à la théorie des schémas de Hilbert, nous permettra
par exemple de construire les schémas de Picard (définis dans [2], II,
n\textsuperscript{o} 3) sous des conditions assez générales.

En résumé, on peut dire qu'on dispose maintenant d'une technique de
constructions projectives à peu près satisfaisante, sauf qu'il manque
encore\textnormal{(*)}\begingroup\renewcommand{\thefootnote}{(*)}\stepcounter{footnote}\footnotetext{Voir l'additif à la fin de l'exposé.}\endgroup un théorème de passage au
quotient par des groupes tels que le groupe projectif, opérant "sans points
fixes", (cf. [2], III, n\textsuperscript{o} 8). La situation semble même un peu
meilleure en géométrie analytique (lorsqu'on s'y borne à l'étude des espaces
analytiques "projectifs" sur un espace analytique donné), car pour les
espaces analytiques la difficulté de passage au quotient par un groupe opérant
joliment disparaît. Par ailleurs, en géométrie algébrique comme en géométrie
analytique, il resterait à mettre au point une technique de construction
valable sans hypothèses projectives.

\fsection{1}{\normalfont\uline{Familles limitées de faisceaux : propriétés de permanence}.}

Soient $\symup{k}$ un corps, et $\symup{X}$ un $\symup{k}$-préschéma, de type fini pour simplifier.
Pour toute extension $\symup{K/k}$, on en déduit donc un $\symup{K}$-préschéma
$\symup{X_K=X\otimes_k K}$. Si $\symup{\mathcal F}$ est un faisceau cohérent sur $\symup{X_K}$, et si
$\symup{K'}$ est une extension de $\symup{K}$, alors
$\symup{\mathcal F\otimes_K K'={\mathcal F}_{K'}}$ est un faisceau quasi-cohérent sur
$\symup{X_K\otimes_K K'=X_{K'}}$. Ceci dit, si $\symup{K}$ et $\symup{K'}$ sont deux extensions
quelconques de $\symup{k}$, $\symup{\mathcal F}$ un faisceau quasi-cohérent sur $\symup{X_K}$ et
$\symup{\mathcal F'}$
\src{250}
un faisceau quasi-cohérent sur $\symup{X_{K'}}$, on dira que $\symup{\mathcal F}$ et
$\symup{\mathcal F'}$ sont \uline{équivalents} s'il existe des $\symup{k}$-homomorphismes de
$\symup{K}$, $\symup{K'}$ dans une même extension $\symup{K''}$ de $\symup{k}$, tels que
$\symup{\mathcal F_{K''}}$ et $\symup{\mathcal F'_{K''}}$ soient isomorphes sur $\symup{X_{K''}}$.
C'est là une relation d'équivalence, et on s'intéressera à des classes
d'équivalence de faisceaux pour cette relation, et à des ensembles de classes
d'équivalence. Notons que si $\symup{X}$ est de type fini sur
$\symup{k}$, alors toute classe de faisceaux cohérents peut être définie par un
faisceau cohérent sur un $\symup{X_K}$, où $\symup{K}$ est une extension \uline{de type fini} de $\symup{k}$.
On peut alors, dans la définition des classes de faisceaux cohérents, se borner
aux extensions \uline{algébriquement closes} de $\symup{k}$, et on peut aussi se limiter à une
extension algébriquement close fixée $\symup{\Omega}$ de $\symup{k}$, de degré de
transcendance infini, deux faisceaux cohérents $\symup{\mathcal F}$ et $\symup{\mathcal F'}$
sur $\symup{X_\Omega}$ étant équivalents si et seulement s'il existe un
$\symup{k}$-automorphisme $\symup{\sigma}$ de $\symup{\Omega}$, tel que
$\symup{\mathcal F\otimes_\Omega(\Omega,\sigma)}$ soit isomorphe à $\symup{\mathcal F'}$.
On constate qu'il y a correspondance biunivoque entre les classes de faisceaux
cohérents pour l'une ou l'autre définition.

Soient $\symup{E}$, $\symup{E'}$ deux ensembles de classes de faisceaux cohérents sur $\symup{X}$.
Considérons les classes de tous les faisceaux de la forme
$\symup{\mathcal F\otimes\mathcal F'}$, où $\symup{\mathcal F}$ et $\symup{\mathcal F'}$ sont des
faisceaux cohérents sur un \uline{même} $\symup{X_K}$, la classe de $\symup{\mathcal F}$ étant dans
$\symup{E}$ et celle de $\symup{\mathcal F'}$ dans $\symup{E'}$. On trouve ainsi un ensemble de
classes de faisceaux cohérents, qu'on notera $\symup{E\otimes E'}$. On définit de même
les $\symup{\mathop{\text{\uwave{Tor}}}\nolimits_i(E,E')}$, etc. De façon générale, à toute fonction
$\symup{\mathcal U}$, associant à toute suite $\symup{\mathcal F_1,\cdots,\mathcal F_n}$ de $\symup{n}$
faisceaux cohérents sur un même $\symup{X_K}$ un ensemble
$\symup{\mathcal U(\mathcal F_1,\cdots,\mathcal F_n)}$ de faisceaux cohérents sur $\symup{X_K}$, et
ayant une propriété évidente de compatibilité avec les isomorphismes de
faisceaux et les images inverses par changement de base, on associe une
fonction, notée par le même symbole $\symup{\mathcal U}$, associant à toute suite de $\symup{n}$
ensembles $\symup{E_1,\cdots,E_n}$ de classes de faisceaux cohérents, un ensemble
$\symup{\mathcal U(E_1,\cdots,E_n)}$ de classes de faisceaux cohérents.

Notre but dans ce numéro est de donner une définition de certains ensembles de
classes de faisceaux, qualifiés de \uline{limités}, et de montrer que les
opérations $\symup{\mathcal U}$ les plus courantes, appliquées à des ensembles limités, donnent
encore des ensembles limités.

Soit $\symup{X}$ un préschéma de type fini sur $\symup{S}$ \uline{noethérien}. Pour tout $\symup{s\in S}$, la
fibre $\symup{X_s}$ est un préschéma de type fini sur $\symup{k(s)}$, et nous considérerons des
classes de faisceaux cohérents sur $\symup{X_s}$, au sens précédent. Cela donne un sens
à la locution : classe de faisceaux cohérents sur une fibre de $\symup{X/S}$, et les
locutions analogues. De même, procédant séparément sur chaque fibre, on peut
encore considérer des opérations telles que $\symup{E\otimes E'}$ etc. faisant
correspondre à des systèmes d'ensembles de
\src{251}
classes de faisceaux cohérents sur les fibres de $\symup{X/S}$, un autre ensemble de
classes de faisceaux cohérents sur les fibres de $\symup{X/S}$.

\result{DÉFINITION}{1.1} \uline{Soit} $\symup{E}$ \uline{un ensemble de classes de faisceaux cohérents
sur les fibres de} $\symup{X/S}$. \uline{On dit que} $\symup{E}$ \uline{est limité, s'il existe un
préschéma} $\symup{S'}$ \uline{de type fini sur} $\symup{S}$, \uline{et un faisceau cohérent} $\symup{\mathcal F'}$ \uline{sur}
$\symup{X'=X\times_S S'}$, \uline{tel que} $\symup{E}$ \uline{soit contenu dans l'ensemble des classes de
faisceaux sur les fibres de} $\symup{X/S}$ \uline{défini par} $\symup{\mathcal F'}$.

Ce dernier, par définition, fait correspondre à un $\symup{s\in S}$ les classes des
faisceaux $\symup{\mathcal F'\otimes_{S'}k(s')}$, où $\symup{s'}$ parcourt les points de $\symup{S'}$
au-dessus de $\symup{s}$ (de sorte que $\symup{k(s')}$ est une extension de $\symup{k(s)}$, et
$\symup{X_s\otimes_{k(s)}k(s')}$ s'identifie à la fibre
$\symup{X'\otimes_{S'}k(s')=X'_{s'}}$ de $\symup{X'}$ en $\symup{s'}$). On peut dire que les familles
limitées sont celles qui sont contenues dans une \uline{famille algébrique} de
faisceaux cohérents, paramétrée par un $\symup{S'}$ de type fini sur $\symup{S}$.

Une réunion finie de familles limitées est limitée (prendre le préschéma somme
des préschémas de paramètres $\symup{S_i}$ définissant les familles algébriques
majorantes). Un changement de base $\symup{T\to S}$ transforme une famille
limitée relative à $\symup{X/S}$ en une famille limitée relative à $\symup{X_T/T}$, et la
réciproque est vraie si $\symup{T\to S}$ est surjectif (ou plus
généralement, si son image contient les $\symup{s}$ qui interviennent effectivement
dans la famille donnée $\symup{E}$ pour $\symup{X/S}$). Cela ramène théoriquement la
détermination des familles limitées au cas où $\symup{S}$ est le spectre d'une algèbre
de type fini sur l'anneau des entiers $\symup{Z}$.

Si $\symup{E}$ et $\symup{E'}$ sont des familles limitées de classes de faisceaux relatives à
$\symup{X/S}$, alors $\symup{E\otimes E'}$ est également limité : en effet, si $\symup{E}$ et $\symup{E'}$ sont
majorés respectivement par les familles algébriques définies par des
$\symup{T\to S}$ et $\symup{F}$ sur $\symup{X_T}$, $\symup{T'\to S}$ et
$\symup{F'}$ sur $\symup{X_{T'}}$, on voit que $\symup{E\otimes E'}$ est majoré par la famille
algébrique définie par $\symup{T''=T\times_S T'\to S}$ et le faisceau
$\symup{F''}$ sur $\symup{X_{T''}}$ produit tensoriel des images inverses de
$\symup{\mathcal F}$ et $\symup{\mathcal F'}$ sur $\symup{X_{T''}}$. Ce raisonnement n'est correct que
parce que le foncteur $\symup{\mathcal F\otimes\mathcal F'}$ en $\symup{\mathcal F}$,
$\symup{\mathcal F'}$ est exact à droite, donc commute à l'extension de la base (et en
particulier au passage aux fibres). Il n'est pas applicable tel quel aux
opérations locales telles que $\symup{\mathop{\text{\uwave{Tor}}}\nolimits_i(E,E')}$,
$\symup{\mathop{\text{\uwave{Hom}}}\nolimits(E,E')}$, $\symup{\mathop{\text{\uwave{Ext}}}\nolimits^i(E,E')}$. On peut cependant
montrer que ces opérations transforment encore des ensembles limités en ensembles
limités, en procédant comme pour $\symup{E\otimes E'}$, mais en utilisant en plus des
résultats du type suivant (tous contenus dans [3], IV, 6.11) : une famille
limitée $\symup{E}$ est toujours majorée par une famille algébrique définie par un
faisceau cohérent $\symup{\mathcal F}$ sur un $\symup{X_T}$ ($\symup{T}$ de type fini sur $\symup{S}$) qui est
\uline{plat} relativement à $\symup{T}$. (On "coupe en morceaux" l'espace des paramètres
initial). De telles
\src{252}
propriétés de platitude sur des faisceaux convenables assurent en effet la
commutation d'opérations telles que $\symup{\mathop{\text{\uwave{Tor}}}\nolimits_i(\mathcal F,\mathcal
F')}$ au changement de base quelconque. La même méthode s'applique pour des
opérations de nature globale : images directes et images directes dérivées de
faisceaux cohérents par des morphismes propres, $\symup{\operatorname{Ext}}$ globaux
relatifs à des morphismes propres (cf. [5], III, \S~6), etc., toutes ces
opérations transforment des familles limitées de faisceaux en familles limitées
de faisceaux (N. B. - Ici les préschémas sur lesquels on prend les divers
faisceaux peuvent changer par les opérations du type envisagé).

Les deux énoncés qui suivent se démontrent essentiellement par la même
technique de platitude ; pour la décomposition primaire sur les fibres d'un
morphisme de type fini, voir en particulier [5], IV.

\result{PROPOSITION}{1.2} \uline{Soient} $\symup{E}$, $\symup{E'}$ \uline{des ensembles limités de classes de
faisceaux sur les fibres de} $\symup{X/S}$, $\symup{X}$ \uline{étant supposé propre sur} $\symup{S}$. \uline{Alors} :

\medskip
\noindent (i) \uline{La famille des noyaux, conoyaux, images d'homomorphismes}
$\symup{\mathcal F\to\mathcal F'}$, \uline{où la classe de} $\symup{\mathcal F}$ \uline{est dans}
$\symup{E}$ \uline{et celle de} $\symup{\mathcal F'}$ \uline{dans} $\symup{E'}$, \uline{est limitée}.

\smallskip
\noindent (ii) \uline{La famille des faisceaux} $\symup{\mathcal F''}$ \uline{extensions d'un}
$\symup{\mathcal F}$ \uline{par un} $\symup{\mathcal F'}$, \uline{où la classe de} $\symup{\mathcal F}$ \uline{est dans} $\symup{E}$
\uline{et celle de} $\symup{\mathcal F'}$ \uline{dans} $\symup{E'}$, \uline{est limitée}.

Moyennant un changement de base convenable, on peut supposer $\symup{E}$ et $\symup{E'}$
définis respectivement par un faisceau cohérent $\symup{\mathcal G}$ et $\symup{\mathcal G'}$
sur un $\symup{X_T/T}$, $\symup{T}$ de type fini sur $\symup{S}$. De plus, on peut supposer satisfaites
certaines hypothèses de platitude, impliquant que la formation du faisceau
$\symup{f_T(\mathop{\text{\uwave{Hom}}}\nolimits_{O_{X_T}}(\mathcal G,\mathcal G'))}$ et
$\symup{\mathop{\text{\uwave{Ext}}}\nolimits^{1}_{f_T}(\mathcal G,\mathcal G')}$ commute au changement
de base pour un morphisme quelconque $\symup{T'\to T}$. De plus, on peut
supposer les faisceaux cohérents précédents sur $\symup{T}$ localement libres. Soient
alors $\symup{T_0}$ et $\symup{T_1}$ les fibrés vectoriels sur $\symup{T}$ dont les faisceaux de germes
de sections sont respectivement les faisceaux précédents. On définit alors
canoniquement un homomorphisme
$\symup{\mathcal G_{T_0}\to\mathcal G'_{T_0}}$ de faisceaux cohérents sur
$\symup{X_{T_0}}$, et une extension
\[
\symup{0\to\mathcal G'_{T_1}\to\mathcal G''
\to\mathcal G_{T_1}\to0}
\]
de faisceaux cohérents sur $\symup{X_{T_1}}$, ayant une propriété universelle évidente.
Ce deuxième faisceau définit une famille algébrique qui majore la famille
envisagée dans (ii). Il en est de même pour les noyau, conoyau et image de
l'homomorphisme
\src{253}
précédent, dont la considération établit (i) (pourvu qu'on suppose le conoyau
plat relativement à $\symup{T_0}$, auquel cas on se ramène encore en découpant $\symup{T_0}$ en
morceaux\ldots).

\result{PROPOSITION}{1.3} \uline{Soit} $\symup{E}$ \uline{une famille limitée de classes de faisceaux
sur les fibres de} $\symup{X/S}$. \uline{Alors les classes des faisceaux structuraux des}
$\symup{(\operatorname{supp}\mathcal F)_{\mathrm{réd}}}$, \uline{où} $\symup{\mathcal F}$ \uline{est un
faisceau cohérent sur un} $\symup{X_K}$, \uline{avec} $\symup{K}$ \uline{algébriquement clos, dont la classe est
dans} $\symup{E}$, \uline{forment une famille limitée}.

Ici, $\symup{(\operatorname{supp}\mathcal F)_{\mathrm{réd}}}$ désigne le support de
$\symup{\mathcal F}$, muni de la structure réduite induite, i. e. son faisceau
structural est le quotient de $\symup{\mathcal O_{X_K}}$ par le plus grand faisceau
d'idéaux définissant $\symup{\operatorname{supp}\mathcal F}$. On peut prouver le
résultat analogue à (1.3) pour les faisceaux déduits canoniquement de
$\symup{\mathcal F}$ par la théorie de la décomposition primaire, par exemple les
$\symup{\mathcal F/\mathcal F_i}$, où les $\symup{\mathcal F_i}$ sont les sous-faisceaux de
$\symup{\mathcal F}$ primaires pour les composantes du support de $\symup{\mathcal F}$,
minimaux pour cette propriété ; ou les $\symup{\mathcal O_{X_K}/\text{\uwave{p}}}$, où
$\symup{\text{\uwave{p}}}$ est un faisceau d'idéaux premier associé à $\symup{\mathcal F}$, ou les
$\symup{\mathcal O_{X_K}/\text{\uwave{q}}}$, où $\symup{\text{\uwave{q}}}$ est un faisceau d'idéaux
primaire associé à une composante du support de $\symup{\mathcal F}$, (le corps de
référence étant algébriquement clos).

\fsection{2}{\normalfont\uline{Familles limitées et polynôme de Hilbert}.}

Dans la suite, nous supposons que $\symup{X}$ est projectif sur $\symup{S}$, muni d'un faisceau
très ample, noté $\symup{\mathcal O_X(1)}$. Pour toute extension $\symup{K}$ d'une extension
résiduelle $\symup{k(s)}$ d'un point $\symup{s}$ de $\symup{S}$, on considère sur $\symup{X_K}$ le faisceau
$\symup{\mathcal O_{X_K}(1)}$ correspondant, qui sera encore très ample.

A tout faisceau cohérent $\symup{\mathcal F}$ sur $\symup{X_K}$, on associe la fonction
\[
\symup{P_{\mathcal F}(n)}=\text{caractéristique d'Euler-Poincaré de $\symup{F(n)}$
sur $\symup{X_K}$}
\]
qui est un polynôme en l'entier $\symup{n}$, appelé \uline{polynôme de Hilbert} de
$\symup{\mathcal F}$. Pour les grandes valeurs de $\symup{n}$, $\symup{P(n)}$ n'est autre chose que la
dimension de $\symup{H^0(X_K,\mathcal F(n))}$ sur $\symup{K}$, puisque les
$\symup{H^i(X_K,\mathcal F(n))}$ sont nuls pour $\symup{i>0}$ et $\symup{n}$ grand.

Si maintenant $\symup{\mathcal F}$ est un faisceau cohérent sur $\symup{X}$, \uline{plat} par rapport
à $\symup{S}$, alors les polynômes de Hilbert des faisceaux $\symup{\mathcal F_s}$ induits sur
les fibres $\symup{X_s}$ relatifs à une même composante connexe de $\symup{S}$, sont tous égaux
[5], III, \S~7. Il en résulte (sans hypothèse de platitude) que l'ensemble des
polynômes de Hilbert des faisceaux $\symup{\mathcal F_s}$, $\symup{s\in S}$, est fini pour tout
faisceau cohérent $\symup{\mathcal F}$ sur $\symup{X}$.
\src{254}
Rappelons d'autre part que si $\symup{\mathcal F}$ est un faisceau cohérent sur $\symup{X}$, il
est isomorphe à un faisceau quotient d'un faisceau $\symup{\mathcal O_X(-n)^N}$, pour
$\symup{n}$, $\symup{N}$ assez grands. Donc les faisceaux $\symup{\mathcal F_s}$ induits sur les fibres
sont également quotients du faisceau $\symup{\mathcal O(-n)^N}$ sur la fibre.

De ces deux remarques, on déduit la partie "il faut" du théorème suivant :

\result{THÉORÈME}{2.1} \uline{Soit} $\symup{X}$ \uline{projectif sur} $\symup{S}$ \uline{noethérien},
$\symup{\mathcal O_X(1)}$ \uline{très ample sur} $\symup{X}$ \uline{relativement à} $\symup{S}$. \uline{Soit} $\symup{E}$ \uline{un ensemble
de classes de faisceaux sur les fibres de} $\symup{X/S}$. \uline{Pour que} $\symup{E}$ \uline{soit limité, il
faut qu'il
satisfasse les conditions suivantes} :

\medskip
\noindent (a) \uline{Il existe un faisceau cohérent} $\symup{\mathcal L}$ \uline{sur} $\symup{X}$ (\uline{qu'on peut
supposer de la forme} $\symup{\mathcal O_X(-n)^N}$) \uline{tel que} $\symup{E}$ \uline{soit contenu dans la
famille des classes de faisceaux cohérents quotients de faisceaux de la forme}
$\symup{\mathcal L_K}$.

\smallskip
\noindent (b) \uline{Les polynômes de Hilbert} $\symup{P_{\mathcal F}}$ \uline{des faisceaux}
$\symup{\mathcal F}$ \uline{dont la classe est dans} $\symup{E}$, \uline{sont éléments d'un même ensemble} fini
\uline{de polynômes}.

Il reste à prouver le "il suffit", qui sera un cas particulier d'un résultat
plus précis. Pour tout Module cohérent $\symup{\mathcal F}$ sur un préschéma de type
fini sur un corps $\symup{K}$, et tout entier $\symup{r}$, soit $\symup{N_r}$ le sous-Module de
$\symup{\mathcal F}$ dont les sections sur un ouvert sont les sections de $\symup{\mathcal F}$
sur cet ouvert dont le support est de dimension $\symup{<r}$. On a donc
$\symup{N_r=\mathcal F}$ pour $\symup{r>\dim\operatorname{supp}\mathcal F}$, $\symup{N_r=0}$ pour
$\symup{r\leq0}$, et on obtient une filtration croissante finie de $\symup{\mathcal F}$ dont
les facteurs $\symup{N_{r+1}/N_r}$ ont comme cycles premiers associés exactement les cycles premiers
associés à $\symup{\mathcal F}$ qui sont de dimension $\symup{r}$. On posera
\[
\symup{\mathcal F_{(r)}=\mathcal F/N_r,}
\]
donc $\symup{\mathcal F_{(r)}}$ a comme cycles premiers associés exactement les cycles
premiers associés à $\symup{\mathcal F}$ qui sont de dimension $\symup{\geq r}$, et en
particulier, il est égal à $\symup{\mathcal F}$ si et seulement si les cycles premiers
associés à $\symup{\mathcal F}$ sont de dimension $\symup{\geq r}$. Ceci posé :

\result{THÉORÈME}{2.2} \uline{Sous les conditions de} \uline{(2.1)}, \uline{soit} $\symup{s}$ \uline{un entier, et
supposons que} $\symup{E}$ \uline{satisfasse les conditions (a), et la forme affaiblie suivante
de (b}) :

\medskip
\noindent $\symup{(b_s)}$ \uline{Les polynômes de Poincaré} $\symup{P_{\mathcal F}}$ \uline{des faisceaux}
$\symup{\mathcal F}$ \uline{dont la classe est dans} $\symup{E}$, \uline{ont des coefficients} en degrés
$\symup{\geq s-1}$ \uline{qui restent bornés}.

\uline{Sous ces conditions, les faisceaux} $\symup{\mathcal F_{(s)}}$ (\uline{lorsque la classe de}
$\symup{\mathcal F}$ \uline{reste dans} $\symup{E}$) \uline{forment une famille limitée. De plus, les
coefficients en degré} $\symup{s-2}$ \uline{des} $\symup{P_{\mathcal F}}$ \uline{restent minorés}.
\src{255}
Ainsi :

\result{COROLLAIRE}{2.3} \uline{Supposons que les faisceaux} $\symup{\mathcal F}$ \uline{dont la classe
est dans} $\symup{E}$ \uline{soient tels que tous les cycles premiers associés soient de
dimension} $\symup{d}$ \uline{satisfaisant} $\symup{s\leq d\leq r}$. \uline{Alors dans} \uline{(2.1)} (\uline{b), on peut se
borner aux coefficients des} $\symup{P_{\mathcal F}}$ \uline{en degrés compris entre} $\symup{s-1}$ \uline{et}
$\symup{r}$.

La fin du numéro est consacrée à l'esquisse de la démonstration de (2.2). Les
lemmes-clefs sont les deux lemmes suivants, dont le premier est bien connu (et
résume le contenu mathématique utile des coordonnées de Chow) :

\result{LEMME 2.4 (CHOW)}{} \uline{Considérons les faisceaux structuraux de
sous-schémas} $\symup{Y}$ \uline{de fibres} $\symup{X_K}$ ($\symup{K}$ \uline{extensions algébriquement closes de corps
résiduels de} $\symup{S}$) \uline{avec} $\symup{Y}$ \uline{réduit et toutes ses composantes de même dimension}
$\symup{r}$, ($\symup{\mathcal O_Y}$ \uline{étant considéré comme un faisceau quotient de}
$\symup{\mathcal O_X}$). \uline{Si les degrés des} $\symup{Y}$ \uline{restent bornés, les} $\symup{Y}$ \uline{forment une
famille limitée}.

Ici, le degré $\symup{a}$ de $\symup{Y}$ peut se définir le plus commodément par le coefficient
du terme dominant de $\symup{P_{\mathcal O_Y}=an^r/r!+\ldots}$.

\result{LEMME}{2.5} \uline{Soient} $\symup{\mathcal L}$ \uline{un faisceau cohérent sur} $\symup{X}$, $\symup{E}$ \uline{un
ensemble de classes de faisceaux quotient} $\symup{\mathcal F}$ \uline{de faisceaux}
$\symup{\mathcal L_K}$ ($\symup{K}$ \uline{extension résiduelle de} $\symup{S}$). \uline{On suppose les fibres de} $\symup{X}$
\uline{sur} $\symup{S}$ \uline{de dimension} $\symup{\leq r}$, \uline{et on pose}
\[
\symup{P_{\mathcal F}(n)=a_{\mathcal F}n^r/r!+b_{\mathcal F}n^{r-1}/(r-1)!}+\text{termes de degré $\symup{<r-1}$}.
\]
\uline{Alors le coefficient} $\symup{a_{\mathcal F}}$ \uline{reste borné, et} $\symup{b_{\mathcal F}}$ \uline{reste
minoré. Si} $\symup{b_{\mathcal F}}$ \uline{reste borné, alors la famille des}
$\symup{\mathcal F_{(r)}}$ \uline{est limitée}.

DÉMONSTRATION. - Quitte à remplacer $\symup{S}$ par une réunion de sous-schémas
de $\symup{S}$ recouvrant $\symup{S}$, on peut supposer qu'il existe un morphisme \uline{fini}
$\symup{f:X\to\mathfrak P^r_S}$, tel que $\symup{\mathcal O_X(1)}$ soit isomorphe à
l'image inverse de $\symup{\mathcal O_{\mathfrak P^r_S}(1)}$, donc pour tout faisceau
cohérent $\symup{\mathcal F}$ sur $\symup{X}$, on a
$\symup{P_{\mathcal F}=P_{f_*(\mathcal F)}}$. D'autre part, on vérifie facilement (par
la technique du numéro précédent) qu'un ensemble de faisceaux $\symup{\mathcal F}$ sur
$\symup{X}$ est limité si et seulement si l'ensemble des $\symup{f_*(\mathcal F)}$ l'est.
Enfin on a
\[
\symup{f_*(\mathcal F)_{(r)}=f_*(\mathcal F_{(r)}).}
\]
\src{256}
Cela nous ramène donc au cas où $\symup{X=\mathfrak P^r_S}$. De plus, on peut supposer
que $\symup{\mathcal L=\mathcal O_{\mathfrak P^r_S}(k)^s}$ pour $\symup{k}$, $\symup{s}$ convenables. Le
coefficient $\symup{a_{\mathcal F}}$ satisfait
\[
\symup{0\leq a_{\mathcal F}\leq s}
\]
et reste donc borné. Cela dit, il revient au même de dire que les
$\symup{P_{\mathcal F}(n)}$ ont un coefficient de $\symup{n^{r-1}}$ restant minoré (resp.
borné) ou de le dire pour les
$\symup{P_{\mathcal F}(n-k)=P_{\mathcal F(-k)}(n)}$. Cela nous ramène au cas où
\[
\symup{\mathcal L=\mathcal O_{\mathfrak P^r_S}^{s}.}
\]

Considérons la suite exacte
\[
\symup{0\to N_r\to\mathcal F
\to\mathcal F_{(r)}\to0,}
\]
d'où
\[
\symup{P_{\mathcal F}=P_{\mathcal F_{(r)}}+P_{N_r},}
\]
et comme le coefficient de $\symup{n^{r-1}}$ dans $\symup{P_{N_r}}$ est positif (car
$\symup{\dim\operatorname{supp}N_r\leq r-1}$), on a
\[
\symup{b_{\mathcal F_{(r)}}\leq b_{\mathcal F}.}
\]
Cela nous permet, pour prouver les assertions du lemme, de remplacer
$\symup{\mathcal F}$ par $\symup{\mathcal F_{(r)}}$, i. e. de supposer que les quotients
$\symup{\mathcal F}$ envisagés de $\symup{\mathcal L}$ sont sans torsion.

Comme $\symup{\mathfrak P^r_K}$ est normal, il s'ensuit que $\symup{\mathcal F}$ est localement
libre de rang $\symup{a=a_{\mathcal F}}$ dans un ouvert
$\symup{U=\mathfrak P^r_K-Y}$, où $\symup{Y}$ est de codimension $\symup{\geq2}$. Donc
$\symup{\mathop{\bigwedge}\limits^{a}\mathcal F}$ est un faisceau sur $\symup{P^r_K}$ dont la restriction
à $\symup{U}$ est inversible, donc ($\symup{\mathfrak P^r_K}$ étant régulier et $\symup{Y}$ de
codimension $\symup{\geq2}$) isomorphe à la restriction d'un faisceau inversible sur
$\symup{\mathfrak P^r_K}$, défini à un isomorphisme près. Ce dernier est de la forme
$\symup{\mathcal O_{\mathfrak P^r_K}(d)}$ pour un entier $\symup{d}$ bien déterminé. Comme
$\symup{\mathop{\bigwedge}\limits^{a}\mathcal F}$ est un quotient de
$\symup{\mathop{\bigwedge}\limits^{a}\mathcal O_{\mathfrak P^r_K}^{s}\simeq
\mathcal O_{\mathfrak P^r_K}^{N}}$, avec $\symup{N=\binom{s}{a}}$, il admet $\symup{N}$ sections
canoniques, définissant donc des sections de
\src{257}
$\symup{\mathcal O_{\mathfrak P^r_K}(d)}$ sur $\symup{U}$, qui sont restrictions de sections
$\symup{s_i}$ ($\symup{1\leq i\leq N}$) de $\symup{\mathcal O_{\mathfrak P^r_K}(d)}$, (puisque
$\symup{\mathfrak P^r_K}$ est normal et $\symup{Y}$ de codimension $\symup{\geq2}$). Ces dernières
engendrent $\symup{\mathcal O_{\mathfrak P^r_K}(d)}$ aux points de $\symup{U}$, donc ne sont pas
toutes nulles, ce qui implique que l'on a $\symup{d\geq0}$. D'autre part, un calcul
facile montre qu'on a
\[
\symup{b_{\mathcal F}=a_{\mathcal F}(r+1)/2+d.}
\]
Cela montre en particulier que $\symup{b_{\mathcal F}\geq0}$, donc $\symup{b_{\mathcal F}}$ est
minoré. Il reste borné si et seulement si $\symup{d}$ reste borné, montrons qu'alors
$\symup{\mathcal F}$ reste dans une famille limitée. On peut supposer
$\symup{a_{\mathcal F}}$ et $\symup{b_{\mathcal F}}$ fixés, soient $\symup{a}$ et $\symup{b}$, donc $\symup{d}$ fixé.
La connaissance des $\symup{N}$ sections $\symup{s_i}$ de $\symup{\mathcal O_{\mathfrak P^r_K}(d)}$,
i. e. d'un homomorphisme
$\symup{s:\mathop{\bigwedge}\limits^{a}\mathcal L_K\to\mathcal O_{\mathfrak P^r_K}(d)}$, permet de récupérer $\symup{\mathcal F}$ comme la co-image de l'homomorphisme composé
correspondant :
\[
\symup{\mathcal L_K\to
\mathop{\text{\uwave{Hom}}}\nolimits\!\left(\mathop{\bigwedge}\limits^{a-1}\mathcal L_K,
\mathop{\bigwedge}\limits^{a}\mathcal L_K\right)\to
\mathop{\text{\uwave{Hom}}}\nolimits\!\left(\mathop{\bigwedge}\limits^{a-1}\mathcal L_K,
\mathcal O_{\mathfrak P^r_K}(d)\right),}
\]
où le premier est l'homomorphisme canonique provenant du produit extérieur, et
le deuxième est déduit de $\symup{s}$. On conclut alors par (1.2) (i).

La conjonction des deux lemmes précédents permet de prouver :

\result{LEMME}{2.6} \uline{Supposons sous les conditions préliminaires de} \uline{(2.1)} \uline{que
l'on ait pour tout} $\symup{\mathcal F}$ :
\[
\symup{P_{\mathcal F}(n)=a_{\mathcal F}n^r/r!+b_{\mathcal F}n^{r-1}/(r-1)!}+\text{termes de degré $\symup{<r-1}$},
\]
\uline{et que les coefficients} $\symup{a_{\mathcal F}}$ \uline{restent bornés. Alors les coefficients}
$\symup{b_{\mathcal F}}$ \uline{restent minorés. Si les} $\symup{b_{\mathcal F}}$ \uline{sont bornés, alors
les} $\symup{\mathcal F_{(r)}}$ \uline{sont limités}.

On peut supposer que les corps de base $\symup{K}$ pour les faisceaux $\symup{\mathcal F}$ sont
algébriquement clos. Munissons chaque
$\symup{\operatorname{supp}\mathcal F_{(r)}=}$ réunion des composantes de dimension
$\symup{r}$
de $\symup{\operatorname{supp}\mathcal F}$, de la structure réduite induite. Alors les
degrés des $\symup{\operatorname{supp}\mathcal F_{(r)}}$ restent majorés par $\symup{a}$, donc
en vertu de (2.4) les $\symup{\operatorname{supp}\mathcal F_{(r)}}$ forment un ensemble
limité. De plus, pour toute composante de
$\symup{\operatorname{supp}\mathcal F_{(r)}}$, la longueur de $\symup{\mathcal F_{(r)}}$ pour
cette composante est $\symup{\leq a}$, donc si $\symup{\mathcal J_{\mathcal F}}$ est l'Idéal
qui définit $\symup{\operatorname{supp}\mathcal F_{(r)}}$, alors $\symup{\mathcal F_{(r)}}$
\src{258}
peut être considéré comme un Module sur le sous-schéma $\symup{Y_{\mathcal F}}$ de $\symup{X}$
défini par $\symup{\mathcal J_{\mathcal F}^{a}}$. Comme dans le lemme précédent, on se
ramène aussi au cas où on a $\symup{\mathcal F=\mathcal F_{(r)}}$, donc $\symup{\mathcal F}$
provient d'un Module sur $\symup{Y_{\mathcal F}}$. Les $\symup{Y_{\mathcal F}}$ correspondent à
une famille limitée de Modules quotients des $\symup{\mathcal O_{X_K}}$, donc
proviennent d'un sous-schéma fermé $\symup{Y}$ d'un schéma $\symup{X\times_S T}$. On peut alors
appliquer (2.5) à $\symup{Y/T}$ et $\symup{\mathcal L\otimes_XY}$, d'où la conclusion.

Nous pouvons maintenant prouver (2.2) par récurrence sur la borne supérieure
$\symup{r}$ des $\symup{\dim\operatorname{supp}\mathcal F}$. L'énoncé est trivial pour $\symup{r<0}$,
supposons donc $\symup{r\geq0}$ et l'énoncé prouvé pour les $\symup{r'<r}$. En vertu de (2.6)
les $\symup{\mathcal F_{(r)}}$ forment un ensemble limité, donc aussi en vertu de (1.2)
(i) les noyaux des homomorphismes
$\symup{\mathcal L_K\to\mathcal F_{(r)}}$ ; il existe donc un Module
cohérent $\symup{\mathcal L'}$ sur $\symup{X}$, tel que les noyaux en question, donc aussi les
$\symup{N_r(\mathcal F)=\operatorname{Ker}(\mathcal F\to\mathcal F_{(r)})}$, soient quotients des Modules $\symup{\mathcal L'_K}$. Comme les $\symup{\mathcal F_{(r)}}$ sont
limités, les $\symup{P_{\mathcal F_{(r)}}}$ restent bornés, et la formule
\[
\symup{P_{\mathcal F}=P_{\mathcal F_{(r)}}+P_{N_r}}
\]
montre alors que les $\symup{P_{N_r}}$ satisfont à la même condition $\symup{(b_s)}$ que les
$\symup{P_{\mathcal F}}$. Donc les $\symup{N_r}$ satisfont les conditions (a) et $\symup{(b_s)}$, et
par l'hypothèse de récurrence les $\symup{(N_r)_{(s)}}$ restent limités. Or
$\symup{\mathcal F_{(s)}}$ est une extension de $\symup{\mathcal F_{(r)}}$ par
$\symup{(N_r)_{(s)}}$, donc par (1.2) (ii), les $\symup{\mathcal F_{(s)}}$ restent limités.
Pour la dernière assertion de (2.2), on note que les noyaux $\symup{N_s}$ de
$\symup{\mathcal F\to\mathcal F_{(s)}}$ sont quotients des noyaux des
homomorphismes $\symup{\mathcal L_K\to\mathcal F_{(s)}}$, qui restent limités en vertu de (1.2)
(i), et que le coefficient du terme en $\symup{n^{s-1}}$ dans $\symup{P_{N_s}}$ reste borné ;
le lemme 2.6 prouve alors que le coefficient du terme suivant reste minoré.
Cela achève la démonstration.

\fsection{3}{\normalfont\uline{Schémas de Hilbert : définition, théorème d'existence}.}

Soient $\symup{X}$ un préschéma sur un autre $\symup{S}$, et $\symup{\mathcal F}$ un Module
quasi-cohérent sur $\symup{X}$. Désignons par
\[
\symup{\operatorname{Quot}(\mathcal F/X/S)}
\]
l'ensemble des Modules quasi-cohérents quotients de $\symup{\mathcal F}$ qui sont plats
sur $\symup{S}$. Soit maintenant $\symup{S'\to S}$ un morphisme de changement de
base, posons $\symup{X'=X\times_S S'}$ et
$\symup{\mathcal F'=\mathcal F\otimes_{\mathcal O_S}\mathcal O_{S'}}$, donc $\symup{X'}$ est
un préschéma sur $\symup{S'}$ muni d'un Module
\src{259}
quasi-cohérent $\symup{\mathcal F'}$, et on peut considérer
$\symup{\operatorname{Quot}(\mathcal F'/X'/S')}$. On pose
\[
\symup{\mathop{\text{\uwave{Quot}}}\nolimits_{\mathcal F/X/S}(S')
=\operatorname{Quot}(\mathcal F'/X'/S')}
\qquad\text{(où $\symup{X'=X\times_SS'}$)}.
\]
Si maintenant $\symup{S''\to S'}$ est un $\symup{S}$-morphisme, alors
$\symup{X''=X\times_SS''}$ est isomorphe à $\symup{X'\times_{S'}S''}$ et $\symup{\mathcal F''}$ à
$\symup{\mathcal F'\otimes_{\mathcal O_{S'}}\mathcal O_{S''}}$, et comme le foncteur
image inverse
$\symup{\mathcal G'\rightsquigarrow\mathcal G'\otimes_{\mathcal O_{S'}}\mathcal O_{S''}}$
de la catégorie des Modules quasi-cohérents sur $\symup{X'}$ dans la catégorie des
Modules quasi-cohérents sur $\symup{X''}$ est exact à droite et transforme les Modules
$\symup{S'}$-plats en Modules $\symup{S''}$-Plats, on trouve une application naturelle
\[
\symup{\operatorname{Quot}(\mathcal F'/X'/S')\to
\operatorname{Quot}(\mathcal F''/X''/S''),}
\]
donc $\symup{\mathop{\text{\uwave{Quot}}}\nolimits_{\mathcal F/X/S}(S')}$ est un \uline{foncteur
contravariant en} $\symup{S'}$ (préschéma au dessus de $\symup{S}$), à valeurs dans la catégorie
des ensembles. Par la suite, nous supposerons que $\symup{X}$ est \uline{projectif} sur $\symup{S}$
\uline{noethérien}, $\symup{\mathcal F}$ \uline{cohérent}, et nous nous limiterons pour simplifier à
des $\symup{S'}$ sur $\symup{S}$ qui sont localement noethériens.

\result{THÉORÈME}{3.1} \uline{Sous ces conditions, le foncteur contravariant}
$\symup{\mathop{\text{\uwave{Quot}}}\nolimits_{\mathcal F/X/S}}$ \uline{sur la catégorie des}
$\symup{S}$\uline{-préschémas localement noethériens est} représentable \uline{par un} $\symup{S}$\uline{-préschéma}
$\symup{\underline{\operatorname{Quot}}_{\mathcal F/X/S}}$, \uline{somme d'une suite de} $\symup{S}$\uline{-schémas
projectifs (a fortiori} $\symup{\underline{\operatorname{Quot}}_{\mathcal F/X/S}}$ \uline{est localement de
type fini sur} $\symup{S}$).

Nous obtiendrons une telle décomposition de la façon suivante. Soit
$\symup{\mathcal O_X(1)}$ un faisceau inversible sur $\symup{X}$ très ample relativement à
$\symup{S}$. Pour tout polynôme $\symup{P(n)}$ à coefficients rationnels, soit
$\symup{\operatorname{Quot}^P(\mathcal F/X/S)}$ la partie de
$\symup{\operatorname{Quot}(\mathcal F/X/S)}$ formée des quotients cohérents
$\symup{\mathcal G}$ de $\symup{\mathcal F}$ qui sont plats sur $\symup{S}$ et dont le polynôme de
Hilbert en tout $\symup{s\in S}$ est égal à $\symup{P}$. On posera alors
\[
\symup{\mathop{\text{\uwave{Quot}}}\nolimits^P_{\mathcal F/X/S}(S')
=\operatorname{Quot}^P(\mathcal F'/X'/S')}
\]
et on obtient ainsi un sous-foncteur de
$\symup{\mathop{\text{\uwave{Quot}}}\nolimits_{\mathcal F/X/S}}$. La propriété d'invariance
des polynômes de Hilbert rappelée dans le numéro 2 implique ceci : \uline{Pour que}
$\symup{\mathop{\text{\uwave{Quot}}}\nolimits_{\mathcal F/X/S}}$ \uline{soit représentable, il faut
et il suffit que les}
$\symup{\mathop{\text{\uwave{Quot}}}\nolimits^P_{\mathcal F/X/S}}$ \uline{le soient, et alors le}
$\symup{S}$\uline{-préschéma} $\symup{\mathop{\text{\uwave{Quot}}}\nolimits_{\mathcal F/X/S}}$, \uline{qui le représente, est
isomorphe au préschéma somme des}
$\symup{\mathop{\text{\uwave{Quot}}}\nolimits^P_{\mathcal F/X/S}}$ \uline{qui représentent les foncteurs}
$\symup{\mathop{\text{\uwave{Quot}}}\nolimits^P_{\mathcal F/X/S}}$. Ceci posé, (3.1) sera
donc une conséquence du théorème suivant :
\src{260}
\result{THÉORÈME}{3.2} \uline{Avec les notations précédentes, le foncteur}
$\symup{\mathop{\text{\uwave{Quot}}}\nolimits^P_{\mathcal F/X/S}}$ \uline{est représentable par un}
$\symup{S}$\uline{-préschéma projectif} $\symup{\operatorname{Quot}^P_{\mathcal F/X/S}}$.

La suite de ce numéro est consacrée à la démonstration de (3.2).

Soit $\symup{\nu}$ un entier. Pour tout $\symup{S'}$ au-dessus de $\symup{S}$, nous désignons par
$\symup{A_\nu(S')}$ l'ensemble des quotients
$\symup{\mathcal G=\mathcal F'/\mathcal K}$ de
$\symup{\mathcal F'=\mathcal F\otimes_{\mathcal O_S}\mathcal O_{S'}}$, cohérents,
plats sur $\symup{S'}$, de polynôme de Hilbert $\symup{P}$, et satisfaisant les conditions suivantes :
\begin{align*}
\text{a.}\quad &\symup{R^if^\prime_*(\mathcal G(n))=0} &&\text{pour $\symup{i>0}$ et $\symup{n\geq\nu}$}\\
\text{b.}\quad &\symup{R^if^\prime_*(\mathcal K(n))=0} &&\text{pour $\symup{i>0}$ et $\symup{n\geq\nu}$}\\
\text{c.}\quad &\symup{f^\prime_*(\mathcal K(\nu+k))=\mathcal S^\prime_k f^\prime_*(\mathcal K(\nu))} &&\text{pour $\symup{k\geq0}$}.
\end{align*}
Pour cette dernière relation, on suppose que $\symup{X}$ est écrit comme le spectre
premier homogène d'une algèbre graduée quasi-cohérente $\symup{\mathcal S_*}$ sur $\symup{S}$,
à degrés positifs, engendrée par $\symup{\mathcal S_1}$, et on pose
$\symup{\mathcal S' =\mathcal S\otimes_{\mathcal O_S}\mathcal O_{S'}}$, de sorte
que $\symup{X'}$ est le spectre premier homogène de $\symup{\mathcal S'}$. Pour prouver
(3.2), on se ramène d'ailleurs facilement au cas où $\symup{X=\mathfrak P^r_S}$ (puisque
$\symup{S}$ est réunion d'ouverts $\symup{U}$ tels que $\symup{X|U}$ soit un sous-schéma fermé de
$\symup{\mathfrak P^r_U}$ et $\symup{\mathcal O_X(1)}$ induit par
$\symup{\mathcal O_{\mathfrak P^r_U}(1)}$), et où $\symup{\mathcal F}$ est de la forme
$\symup{(\mathcal O_{\mathfrak P^r_S})^{N}}$, donc est plat sur $\symup{S}$. Alors dans ce qui
précède, les faisceaux $\symup{\mathcal K}$ sont également plats sur $\symup{S'}$. Il résulte
alors des relations de Künneth ([5], III, \S~7), et de (b) que les conditions (a) et (b)
sont stables par changement de base, et impliquent que pour $\symup{n\geq\nu}$, la
formation de $\symup{f'_*(\mathcal G(n))}$ et de $\symup{f'_*(\mathcal K(n))}$ commute à
l'extension de la base (loc. cit.), donc (c) est également stable par extension
de la base.

En d'autres termes, $\symup{A_\nu(S')}$ \uline{est un foncteur contravariant en} $\symup{S'}$, de \uline{façon
précise un sous-foncteur de}
$\symup{A(S')=\mathop{\text{\uwave{Quot}}}\nolimits^P_{\mathcal F/X/S}(S')}$. Pour $\symup{\nu}$
variable, on obtient ainsi une suite croissante de parties $\symup{A_\nu(S')}$ de
$\symup{A(S')}$, dont la réunion est $\symup{A(S')}$ en vertu d'un théorème bien connu de
SERRE [5], III, \S~2. Notons maintenant que si $\symup{\mathcal G}$ est un quotient
cohérent de $\symup{\mathcal F'}$, plat sur $\symup{S'}$, $\symup{s}$ un élément de $\symup{S'}$ tel que le
changement de base, $\symup{\operatorname{Spec}(k(s))\to S'}$ donne lieu à
un quotient $\symup{G_s}$ de $\symup{F_s}$ satisfaisant aux conditions (a),
(b), (c), i. e. qui est dans $\symup{A_\nu\operatorname{Spec}(k(s))}$, alors il existe
un voisinage ouvert $\symup{U}$ de $\symup{s}$ tel que ces mêmes conditions soient vérifiées
pour $\symup{\mathcal G|f'^{-1}(U)}$, i. e. ce quotient est dans $\symup{A_\nu(U)}$ ; pour (a)
et (b), cela résulte en effet du "théorème des fonctions holomorphes" [5],
III, \S~7, et (c) résulte du lemme de Nakayama, et du fait qu'on sait en tous
cas que
$\symup{f'_*(\mathcal K(n+k))=\mathcal S'_k f'_*(\mathcal K(n))}$
\src{261}
pour $\symup{n}$ assez grand, et $\symup{k\geq0}$, (cf. [5], III, \S~2). De ces remarques on
conclut ceci (comparer [4], IV). \uline{Pour que le foncteur} $\symup{A}$ \uline{soit représentable,
il faut et il suffit que les foncteurs} $\symup{A_\nu}$ \uline{le soient, et alors le}
$\symup{S}$\uline{-préschéma} $\symup{Q}$ \uline{qui représente} $\symup{A}$ \uline{est réunion croissante d'ouverts} $\symup{Q_\nu}$
\uline{représentant les} $\symup{A_\nu}$.

Soit
\[
\symup{M_* =\sum_{n\geq0}f_*(\mathcal F(n))=\mathcal S_*^{N},}
\]
de sorte que l'on a
\[
\symup{M'_* =M_*\otimes_{\mathcal O_S}\mathcal O_{S'}
=\sum_{n\geq0}f'_*(\mathcal F'(n))=\mathcal S'_*{}^{N}.}
\]
Il résulte de (a) que l'on a :

\noindent (a') $\symup{f'_*(\mathcal G(n))}$ est localement libre de rang $\symup{P(n)}$ pour
$\symup{n\geq\nu}$, [5], III, \S~7, et de (b), pour $\symup{i=1}$ :

\noindent (a'') $\symup{f'_*(\mathcal G(n))}$ est un Module quotient de $\symup{M'_n}$.

D'ailleurs, la connaissance de ce Module quotient, pour $\symup{n=\nu}$, implique en
vertu de (c) celle des sous-Modules $\symup{f'_*(\mathcal K(n))}$ de $\symup{M'_n}$ pour
$\symup{n\geq\nu}$, donc la connaissance de $\symup{\mathcal K}$ et par suite de $\symup{\mathcal G}$.
On obtient ainsi une \uline{application injective}
\[
\symup{A_\nu(S')\to\operatorname{Grass}_{P(\nu)}(M'_\nu)}
\]
de $\symup{A_\nu(S')}$ dans l'ensemble des Modules quotients localement libres de rang
$\symup{P(\nu)}$ de $\symup{M'_\nu}$, d'où un homomorphisme \uline{fonctoriel}
\[
\symup{i_\nu:A_\nu(S')\to
\mathop{\text{\uwave{Grass}}}\nolimits_{P(\nu)}(M_\nu)(S'),}
\]
où le foncteur du deuxième membre est représentable par le schéma grassmanien
$\symup{\operatorname{Grass}_{P(\nu)}(M_\nu)}$ (comparer [4], V), qui est projectif
sur $\symup{S}$. Je dis :

\result{LEMME}{3.3} $\symup{A_\nu(S')}$ \uline{est un foncteur représentable, et le morphisme}
$\symup{Q_\nu\to\operatorname{Grass}_{P(\nu)}(M_\nu)}$ \uline{qui représente
l'homomorphisme} $\symup{i_\nu}$ \uline{est une immersion} (ce qui implique que $\symup{Q_\nu}$ \uline{est
quasi-projectif sur} $\symup{S}$).

Cette assertion est équivalente à la suivante (comparer [4], IV) : Supposons
donné un Module quotient $\symup{N}$ de $\symup{M'_\nu}$, localement libre de rang $\symup{P(\nu)}$,
alors il existe un sous-préschéma $\symup{Z}$ de $\symup{S'}$ tel que pour tout préschéma
localement noethérien $\symup{T'}$
\src{262}
sur $\symup{S'}$, l'image inverse de $\symup{N}$ sur $\symup{T'}$ est dans
$\symup{\operatorname{Im}A_\nu(T')}$ si et seulement si $\symup{T'\to S'}$ est
majoré par le sous-préschéma $\symup{Z}$. Changeant de notations, on peut supposer
$\symup{S'=S}$, i. e. on s'est donné un quotient $\symup{N_\nu}$ de $\symup{M_\nu}$ par un sous-Module
$\symup{R_\nu}$. Pour qu'il provienne d'un élément de $\symup{A_\nu(S)}$, il faut et il suffit
qu'il satisfasse les deux conditions suivantes :

\medskip
\noindent (i) $\symup{M_{\nu+k}/\mathcal S_kR_\nu}$ est localement libre de rang
$\symup{P(\nu+k)}$ pour $\symup{k\geq0}$.

\smallskip
\noindent (ii) Le sous-faisceau $\symup{\mathcal K}$ de $\symup{\mathcal F}$ défini par le
sous-Module gradué
$\symup{R_* =\sum\limits_{k\geq0}\mathcal S_kR_\nu}$
de $\symup{M_*}$ ([5], II, \S~3) et le quotient $\symup{\mathcal G=\mathcal F/\mathcal K}$,
satisfont les conditions (a) et (b) plus haut. Ces conditions sont manifestement nécessaires, d'autre part si elles sont
vérifiées, alors le faisceau $\symup{\mathcal G}$ défini dans (ii), étant isomorphe au
faisceau associé au $\symup{\mathcal S_*}$-Module gradué $\symup{N_*}$ somme des
$\symup{M_{\nu+k}/\mathcal S_kR_\nu}$, est \uline{plat} sur $\symup{S}$ (car ses fibres sont des
facteurs directs de modules localisés de $\symup{N_*}$ pour des idéaux premiers homogènes
de $\symup{\mathcal S_*}$), et correspond au polynôme de Hilbert $\symup{P}$ en vertu de (i).
Compte tenu de (ii), on voit alors que (a') et (a'') sont vérifiées, donc pour
$\symup{n\geq\nu}$, l'homomorphisme naturel
$\symup{N_n\to f_*(\mathcal G(n))}$ est un homomorphisme \uline{surjectif} de
Modules localement libres de même rang, donc un isomorphisme, donc on a
$\symup{f_*(\mathcal K(\nu+k))=\mathcal S_kR_\nu}$ pour tout $\symup{k\geq0}$,
ce qui prouve que $\symup{\mathcal G\in A_\nu(S)}$ et que $\symup{R_\nu}$ est l'élément de
$\symup{\operatorname{Grass}_{P(\nu)}(M_\nu)}$ défini par $\symup{\mathcal G}$.

Ce critère (i) (ii) s'applique également à la situation obtenue après un
changement de base $\symup{S'\to S}$. Nous allons prouver d'abord que le
fait que la condition (i) soit vérifiée après le changement de base
$\symup{S'\to S}$, s'exprime par la condition que $\symup{S'\to S}$
est majoré par un certain sous-préschéma $\symup{Z}$ de $\symup{S}$ ; une fois ce résultat
obtenu, on est ramené (remplaçant $\symup{S}$ par $\symup{Z}$) au cas où la condition (i) est
déjà vérifiée sur $\symup{S}$, et comme elle est stable par changement de base, il
reste à exprimer la condition (ii). Mais alors, si $\symup{U}$ désigne l'ensemble des
$\symup{s\in S}$ tels que la cohomologie des faisceaux induits sur la fibre $\symup{X_s}$ par
$\symup{\mathcal G(n)}$ et $\symup{\mathcal K(n)}$ soit nulle en dimension $\symup{>0}$ pour
$\symup{n\geq\nu}$, on a déjà signalé que $\symup{U}$ est ouvert, et la condition (ii) sera
vérifiée après un changement de base $\symup{S'\to S}$ si et seulement si
$\symup{S'\to S}$ est majoré par $\symup{U}$, ce qui prouve le lemme 3.3. Il reste
donc à prouver le lemme suivant :

\result{LEMME}{3.4} \uline{Soient} $\symup{S}$ \uline{un préschéma localement noethérien, muni d'une
algèbre graduée quasi-cohérente} $\symup{\mathcal S_*}$ \uline{à degrés positifs, engendrée par}
$\symup{\mathcal S_1}$, $\symup{M_*}$ \uline{un} $\symup{\mathcal S}$\uline{-Module gradué quasi-cohérent de type
fini}, $\symup{P}$ \uline{un polynôme à coefficients rationnels}, $\symup{\nu}$ \uline{un entier, Alors il
existe un sous-préschéma} $\symup{Z}$ \uline{de} $\symup{S}$ (\uline{évidemment unique})
\src{263}
\uline{ayant la propriété suivante : Pour tout préschéma} $\symup{S'}$ \uline{sur} $\symup{S}$, \uline{pour que}
$\symup{M_n\otimes_{\mathcal O_S}\mathcal O_{S'}}$ \uline{soit localement libre de rang} $\symup{P(n)}$
\uline{pour tout} $\symup{n\geq\nu}$, \uline{il faut et il suffit que} $\symup{S'\to S}$ \uline{soit majoré
par} $\symup{Z}$.

On peut évidemment supposer $\symup{S}$ affine, donc noethérien. Alors :

\result{LEMME}{3.5} \uline{Pour tout entier} $\symup{N\geq\nu}$, \uline{soit} $\symup{U_N}$ \uline{l'ouvert de} $\symup{S}$
\uline{formé des} $\symup{s\in S}$ \uline{tels que}
$\symup{\operatorname{rang}_{k(s)}M_{ns}\otimes_{\mathcal O_{S,s}}k(s)\leq P(n)}$
\uline{pour tout} $\symup{\nu\leq n\leq N}$. \uline{Alors la suite décroissante des ouverts} $\symup{U_N}$ \uline{est stationnaire}.

On sait [3], IV, que $\symup{S}$ admet une partition finie en sous-préschémas réduits
$\symup{S_i}$, tels que chaque $\symup{M\otimes_{\mathcal O_S}\mathcal O_{S_i}}$ soit \uline{plat} sur
$\symup{S_i}$ : On peut donc supposer $\symup{M}$ plat, donc les $\symup{M_n}$ plats. Enfin, on peut
évidemment supposer $\symup{S}$ connexe. Mais alors il existe un entier $\symup{n_0}$ et un
polynôme $\symup{Q}$ tels que
\[
\symup{\operatorname{rang}_{k(s)}
M_{ns}\otimes_{\mathcal O_{S,s}}k(s)=Q(n)}
\quad\text{pour $\symup{n\geq n_0}$}
\]
[5], III, \S~7. Supposons d'abord $\symup{P\neq Q}$ donc
$\symup{P(n)\neq Q(n)}$ pour $\symup{n}$ grand, alors on a évidemment $\symup{U_N=\varnothing}$ pour $\symup{N}$
grand, a fortiori la suite des $\symup{U_N}$ est stationnaire. Si au contraire $\symup{P=Q}$, alors on aura $\symup{U_N=U_{n_0}}$ pour $\symup{N\geq n_0}$,
et la suite des $\symup{U_N}$ est encore stationnaire.

En particulier, l'ensemble $\symup{U_\infty}$ des $\symup{s\in S}$ tels que
\[
\tag{*}
\symup{\operatorname{rang}_{k(s)}
M_{ns}\otimes_{\mathcal O_{S,s}}k(s)\leq P(n)}
\quad\text{pour tout $\symup{n\geq\nu}$},
\]
étant l'intersection des $\symup{U_N}$, est ouvert. On peut alors, pour la démonstration
de (3.4), remplacer $\symup{S}$ par l'ouvert $\symup{U_\infty}$, ce qui nous ramène au cas où
l'inégalité (*) est satisfaite en tout $\symup{s\in S}$. D'ailleurs :

\result{LEMME}{3.6} \uline{Soit} $\symup{M}$ \uline{un Module sur un préschéma localement noethérien}
$\symup{S}$, \uline{et} $\symup{r}$ \uline{un entier. Alors il existe un sous-préschéma} $\symup{Z}$ \uline{de} $\symup{S}$
(\uline{évidemment unique) ayant la propriété suivante : pour tout} $\symup{S'}$ \uline{sur} $\symup{S}$, \uline{pour
que} $\symup{M\otimes_{\mathcal O_S}\mathcal O_{S'}}$ \uline{soit localement libre de rang} $\symup{r}$,
\uline{il faut et il suffit que} $\symup{S'\to S}$ \uline{soit majoré par} $\symup{Z}$. \uline{Si on a}
$\symup{\operatorname{rang}_{k(s)}M_s\otimes_{\mathcal O_{S,s}}k(s)\leq r}$
\uline{pour tout} $\symup{s}$, \uline{alors} $\symup{Z}$ \uline{est un sous-préschéma fermé de} $\symup{S}$ (on suppose $\symup{M}$ cohérent).
\src{264}
En effet, le raisonnement précédent nous ramène au cas où on a l'inégalité (*)
pour tout $\symup{s\in S}$ (en remplaçant au besoin $\symup{S}$ par la partie ouverte formée
des $\symup{s}$ où l'inégalité est vérifiée). On peut alors supposer que $\symup{M}$ s'insère
dans une suite exacte
\[
\symup{\mathcal O_S^q\to\mathcal O_S^r\to M\to0,}
\]
et la condition envisagée sur les $\symup{S'}$ sur $\symup{S}$ signifie aussi que dans la suite
exacte $\symup{\mathcal O_{S'}^q\to\mathcal O_{S'}^r
\to M'\to0}$ correspondante, la deuxième flèche est un
isomorphisme, i. e. la première est nulle. On voit alors que le sous-préschéma
fermé $\symup{Z}$ de $\symup{S}$ défini par l'Idéal engendré par les coefficients de la matrice
définissant l'homomorphisme
$\symup{\mathcal O_S^q\to\mathcal O_S^r}$, satisfait à la condition voulue.

Revenant alors à la démonstration de (3.4) où nous l'avions laissée, on désigne
par $\symup{Z_n}$ le sous-schéma \uline{fermé} de $\symup{S}$ associé en vertu de (3.6) au Module $\symup{M_n}$
et à l'entier $\symup{r=P(n)}$, par $\symup{Z'_N}$ le $\symup{\operatorname{Inf}}$ des $\symup{Z_n}$ pour
$\symup{\nu\leq n\leq N}$. Alors les $\symup{Z'_N}$ forment une suite décroissante de
sous-schémas fermés de $\symup{S}$, donc nécessairement stationnaire. Soit $\symup{Z}$ la
valeur constante des $\symup{Z'_N}$ pour $\symup{N}$ grand. C'est le $\symup{Z}$ cherché dans (3.4).
Cela achève la démonstration de (3.4), donc aussi celle de (3.3).

On a donc prouvé que
$\symup{\mathop{\text{\uwave{Quot}}}\nolimits^P_{\mathcal F/X/S}}$ \uline{est représentable par un}
$\symup{S}$\uline{-préschéma} $\symup{Q}$ \uline{qui est une réunion croissante de sous-préschémas ouverts}
$\symup{Q_\nu}$ \uline{quasi-projectifs sur} $\symup{S}$. Pour aller plus loin, il faut invoquer le
théorème 2.1, d'où on conclut facilement que $\symup{Q}$ est \uline{quasi-compact}, (car image du
préschéma de type fini $\symup{S'}$ sur $\symup{S}$ qui paramètre la famille des faisceaux
quotients des $\symup{\mathcal F_K}$ ayant le polynôme de Hilbert $\symup{P}$). Donc $\symup{Q}$ est
égal à l'un des $\symup{Q_\nu}$, donc quasi-projectif sur $\symup{S}$. Pour prouver qu'il est
projectif sur $\symup{S}$, il reste donc à prouver qu'il est \uline{propre} sur $\symup{S}$, et pour
cela il suffit d'invoquer le critère valuatif de propreté sous la forme [5],
II, 7.3.8. Il suffit de vérifier ceci :

\result{LEMME}{3.7} \uline{Soient} $\symup{S}$ \uline{le spectre d'un anneau de valuation discrète},
$\symup{s}$ \uline{son point générique}, $\symup{X}$ \uline{un préschéma sur} $\symup{S}$, $\symup{\mathcal F}$ \uline{un Module
quasi-cohérent sur} $\symup{X}$, $\symup{\mathcal G_s}$ \uline{un Module quotient quasi-cohérent de}
$\symup{\mathcal F_s=\mathcal F\otimes_{\mathcal O_S}k(s)}$ \uline{sur} $\symup{X_s}$. \uline{Alors il existe
un Module quotient quasi-cohérent} $\symup{\mathcal G}$ \uline{unique de} $\symup{\mathcal F}$, \uline{plat sur}
$\symup{S}$, \uline{et dont la restriction à} $\symup{X_s}$ \uline{soit} $\symup{\mathcal G_s}$.
\src{265}
En effet, si $\symup{\mathcal G_s=\mathcal F_s/\mathcal K_s}$, il suffit de considérer
le plus grand sous-faisceau $\symup{\mathcal K}$ de $\symup{\mathcal F}$ induisant
$\symup{\mathcal K_s}$, [5], I, 9.4.2 et de prendre
$\symup{\mathcal G=\mathcal F/\mathcal K}$. On vérifie facilement que ce faisceau
convient.

Le théorème 3.2, et par suite (3.1) est complètement démontré.

La démonstration prouve en même temps ceci :

\result{PROPOSITION}{3.8} \uline{Sous les conditions de} \uline{(3.2)}, \uline{soient}
$\symup{Q=\underline{\operatorname{Quot}}^P_{\mathcal F/X/S}}$, $\symup{X_Q=X\times_SQ}$,
$\symup{F_Q=\mathcal F\otimes_{\mathcal O_S}\mathcal O_Q}$, \uline{et soit}
$\symup{\mathcal G}$ \uline{le quotient cohérent de} $\symup{\mathcal F_Q}$, \uline{plat sur} $\symup{Q}$, \uline{ayant le
polynôme de Hilbert relatif} $\symup{P}$, \uline{tel que} $\symup{(Q,\mathcal G)}$ \uline{représente le foncteur}
$\symup{\mathop{\text{\uwave{Quot}}}\nolimits^P_{\mathcal F/X/S}}$. \uline{Il existe un entier}
$\symup{\nu}$ \uline{tel que, pour} $\symup{n\geq\nu}$, $\symup{(f_Q)_*(\mathcal G(n))}$ \uline{soit un Module sur} $\symup{Q}$
\uline{localement libre de rang} $\symup{P(n)}$, \uline{et très ample relativement à} $\symup{S}$, \uline{i. e.
définissant une immersion de} $\symup{Q}$ \uline{dans un schéma grassmanien}
$\symup{\underline{\operatorname{Grass}}_{P(n)}(M)}$ \uline{sur} $\symup{S}$. \uline{A fortiori, pour} $\symup{n\geq\nu}$, \uline{le
faisceau} $\symup{\mathop{\bigwedge}\limits^{P(n)}(f_Q)_*(\mathcal G(n))}$ \uline{sur} $\symup{Q}$ \uline{est inversible très ample relativement à} $\symup{S}$.

On peut en effet se ramener, comme pour (3.2), au cas où $\symup{F}$ est plat
sur $\symup{S}$, et alors il suffit de prendre un entier $\symup{\nu}$ tel que $\symup{A_\nu=A}$ (avec
les notations précédentes).

Le cas le plus important d'application de (3.2) est celui où
$\symup{\mathcal F=\mathcal O_X}$. On écrit alors
\[
\symup{\underline{\operatorname{Quot}}_{\mathcal O_X/X/S}=\underline{\operatorname{Hilb}}_{X/S},
\qquad
\underline{\operatorname{Quot}}^P_{\mathcal O_X/X/S}=\underline{\operatorname{Hilb}}^P_{X/S},}
\]
donc on a une décomposition
\[
\symup{\underline{\operatorname{Hilb}}_{X/S}=\coprod_P\underline{\operatorname{Hilb}}^P_{X/S}.}
\]
Par définition, $\symup{\underline{\operatorname{Hilb}}_{X/S}}$ représente le foncteur
$\symup{\mathop{\text{\uwave{Hilb}}}\nolimits_{X/S}(S')}$ = ensemble des sous-préschémas
fermés de $\symup{X'=X\times_SS'}$ qui sont plats sur $\symup{S'}$ ; et
$\symup{\underline{\operatorname{Hilb}}^P_{X/S}}$ représente le sous-foncteur correspondant aux
sous-préschémas fermés admettant un polynôme de Hilbert donné $\symup{P}$. Ces
préschémas sont aussi appelés le \uline{préschéma de Hilbert} de $\symup{X}$ sur $\symup{S}$,
respectivement le \uline{préschéma de Hilbert d'indice} $\symup{P}$. La terminologie est
justifiée par le rôle joué dans la théorie par les polynômes de Hilbert. Leur
différence de nature avec les classiques variétés de Chow (destinées à
paramétrer des cycles, et non des variétés) est du même ordre qu'entre
\src{266}
\uline{l'anneau de Chow} des classes de cycles d'une variété, et \uline{l'anneau des classes de faisceaux} de la variété (tel qu'il s'introduit dans le théorème de
Riemann-Roch [1]) ; on notera en effet que lorsque $\symup{X=\mathfrak P^r_S}$, $\symup{S}$
étant le spectre d'un corps, la connaissance du polynôme de Hilbert d'un Module
cohérent $\symup{\mathcal F}$ sur $\symup{X}$ équivaut aussi à celle des classes de Chern de
$\symup{\mathcal F}$ et de son rang, ou encore de la classe de $\symup{\mathcal F}$ dans l'anneau des classes
de faisceaux cohérents sur $\symup{X}$.

\result{REMARQUES}{3.9} On notera aussi que la construction des
$\symup{\underline{\operatorname{Quot}}_{\mathcal F/X/S}}$ et
$\symup{\underline{\operatorname{Quot}}^P_{\mathcal F/X/S}}$ a été ramenée au cas où
$\symup{X=\mathfrak P^r_S}$, $\symup{\mathcal F=\mathcal O_X^N}$, $\symup{\mathcal O_X(1)}$ étant le
faisceau très ample habituel ; de façon précise, les
$\symup{\underline{\operatorname{Quot}}_{\mathcal F/X/S}}$ généraux se réalisent comme des
sous-préschémas fermés des précédents. Comme la formation des
$\symup{\underline{\operatorname{Quot}}^P_{\mathcal F/X/S}}$ est évidemment compatible avec les
changements de base $\symup{S'\to S}$, on voit qu'on est ramené au cas où
on a de plus $\symup{S=\operatorname{Spec}(\mathbf Z)}$. En fin de compte, on est donc
ramené à étudier les schémas projectifs sur $\symup{\mathbf Z}$ :
\[
\symup{\mathfrak Q^P_{r,N}=
\underline{\operatorname{Quot}}^{P}_{(\mathcal O_{\mathfrak P^r_{\mathbf Z}})^N/
\mathfrak P^r_{\mathbf Z}/\operatorname{Spec}(\mathbf Z)},}
\]
et plus particulièrement les schémas de Hilbert absolus :
\[
\symup{\underline{\operatorname{Hilb}}^P_r=\mathfrak Q^P_{r,1}.}
\]
Une étude plus détaillée de ces schémas, à commencer par la détermination de
leurs composantes connexes (sont-ils connexes ?), leurs composantes irréductibles (en
vertu de SERRE [7], il peut y avoir des composantes irréductibles qui se
trouvent tout entières sur un nombre premier $\symup{p\neq0}$), serait fort
intéressante. Rappelons la question de WEIL si les composantes irréductibles des
fibres de $\symup{\underline{\operatorname{Hilb}}^P_r}$ sur les
$\symup{s\in\operatorname{Spec}(\mathbf Z)}$ correspondent à des extensions
"régulières" du corps premier, i. e. si elles sont "relativement connexes".
Il n'est pas exclu que ces questions soient plus abordables pour les schémas de
Hilbert que pour les "variétés de Chow".

\fsection{4}{\normalfont\uline{Variantes}.}

\noindent a. Sous les conditions de (3.1), soit $\symup{U}$ un ouvert dans $\symup{X}$, et
désignons par $\symup{A'}$ le sous-foncteur de
$\symup{A=\mathop{\text{\uwave{Quot}}}\nolimits_{\mathcal F/X/S}}$, tel que $\symup{A'(S')}$ soit
l'ensemble des Modules quotients $\symup{\mathcal G}$ de $\symup{\mathcal F'}$, plats sur
$\symup{S'}$, dont le support est contenu dans $\symup{U'}$. On voit aussitôt que $\symup{A'}$ est
représentable par une partie ouverte du préschéma
$\symup{\underline{\operatorname{Quot}}_{\mathcal F/X/S}}$ qui représente $\symup{A}$. Il s'ensuit que les
théorèmes (3.1) et (3.2)
\src{267}
restent valables en supposant que $\symup{X}$ est \\uline{quasi-projectif} sur $\symup{S}$ au lieu de
projectif sur $\symup{S}$, quand on remplace aussi dans les conclusions les mots
"projectif" par "quasi-projectif", et qu'on désigne maintenant par
$\symup{\mathop{\text{\uwave{Quot}}}\nolimits_{\mathcal F/X/S}(S')}$ l'ensemble des
quotients cohérents $\symup{\mathcal G}$ de $\symup{\mathcal F'}$, plats sur $\symup{S'}$, \uline{dont le support est propre sur} $\symup{S'}$.

\medskip
\noindent b. De façon générale, on peut imposer aux quotients $\symup{G}$ de
$\symup{\mathcal F'=\mathcal F\otimes_{\mathcal O_S}\mathcal O_{S'}}$, plats sur $\symup{S'}$
toutes sortes de conditions naturelles supplémentaires, stables par changement
de base, obtenant ainsi autant de sous-foncteurs de
$\symup{\mathop{\text{\uwave{Quot}}}\nolimits_{\mathcal F/X/S}}$ qu'on se propose de
représenter. Le critère habituel permet dans beaucoup de cas de prouver qu'on
obtient encore des foncteurs représentables par des parties ouvertes de
$\symup{\underline{\operatorname{Quot}}_{\mathcal F/X/S}}$. Il en est en particulier ainsi
lorsqu'on impose l'une des propriétés supplémentaires suivantes :

\smallskip
\noindent $\symup{1^\circ}$ Les dimensions des cycles premiers associés aux Modules
$\symup{G_{s'}}$, ($\symup{s'\in S'}$) induits sur les fibres $\symup{X'_{s'}}$ appartiennent
à un ensemble donné d'entiers.

\smallskip
\noindent $\symup{2^\circ}$ (Lorsque $\symup{\mathcal F=\mathcal O_X}$, donc $\symup{\mathcal G}$
correspond à un sous-préschéma fermé $\symup{Y}$ de $\symup{X'}$) ; $\symup{Y}$ est un préschéma \\uline{simple}
([3], IV) sur $\symup{S'}$, resp. \\uline{normal} sur $\symup{S'}$ (i. e. les fibres $\symup{Y_{s'}}$ sont normales
"sur $\symup{k(s')}$", i. e. sont normales par toute extension du corps de base),
resp. (lorsque $\symup{X}$ est plat sur $\symup{S}$) sont des $\symup{S'}$-intersections complètes
locales dans $\symup{X'}$ relativement à $\symup{S'}$ (i. e. les fibres $\symup{Y_{s'}}$ sont des
intersections complètes locales dans les $\symup{X_{s'}}$).

D'autres conditions feraient intervenir des propriétés de nature
cohomologique sur les Modules $\symup{G_{s'}}$ induits sur les $\symup{X'_{s'}}$,
etc. Bien entendu, la conjonction de conditions dont chacune est représentée
par un ouvert $\symup{U_i}$ de $\symup{\underline{\operatorname{Quot}}_{\mathcal F/X/S}}$, est représentée
par l'ouvert intersection. Par exemple, considérant pour tout $\symup{S'}$ sur $\symup{S}$
l'ensemble des sous-préschémas fermés $\symup{Y}$ de $\symup{X'=X\times_SS'}$ qui sont des
revêtements étales ([3], I) de rang donné $\symup{r}$ de $\symup{S'}$, on trouve un foncteur
contravariant représentable en $\symup{S'}$.

\medskip
\noindent c. Les préschémas $\symup{\underline{\operatorname{Hom}}_S(X,Y)}$,
$\symup{\prod\limits_{X/S}(Z/X)}$, $\symup{\underline{\operatorname{Isom}}_S(X,Y)}$, définis dans [2], II, C,
n\textsuperscript{o} 2, existent moyennant des hypothèses projectives
convenables, et se réalisent comme des ouverts dans des préschémas de Hilbert
convenables. Comme on a $\symup{\underline{\operatorname{Hom}}_S(X,Y)=\prod\limits_{X/S}((X\times Y)/X)}$, le cas de $\symup{\underline{\operatorname{Hom}}_S(X,Y)}$ se ramène à celui de $\symup{\prod\limits_{X/S}(Z/X)}$.
On note alors que pour tout $\symup{S'}$ sur $\symup{S}$, l'ensemble des sections de
$\symup{Z'=Z\times_SS'}$ sur $\symup{X'=X\times_SS'}$ est en correspondance biunivoque avec
l'ensemble des sous-préschémas $\symup{\Gamma}$ de $\symup{Z'}$ (nécessairement fermés si $\symup{Z}$
est séparé sur $\symup{X}$) tels que le morphisme $\symup{\Gamma\to X'}$ induit par
$\symup{Z'\to X'}$ soit un isomorphisme. De cette façon,
\src{268}
\uline{lorsque} $\symup{X}$ \uline{est plat et propre sur} $\symup{S}$, \uline{et} $\symup{Z}$ \uline{est quasi-projectif sur} $\symup{S}$,
$\symup{\prod\limits_{X/S}(Z/X)}$ \uline{existe et se réalise comme un sous-préschéma ouvert de}
$\symup{\underline{\operatorname{Hilb}}_{Z/S}}$. \uline{Donc lorsque} $\symup{X}$ \uline{est projectif et plat sur} $\symup{S}$, \uline{et}
$\symup{Y}$ \uline{quasi-projectif sur} $\symup{S}$, $\symup{\underline{\operatorname{Hom}}_S(X,Y)}$ \uline{existe et se réalise
comme un sous-préschéma ouvert de}
$\symup{\underline{\operatorname{Hilb}}_{(X\times_SY)/S}}$. Lorsque $\symup{X}$ et $\symup{Y}$ sont tous les deux
projectifs sur $\symup{S}$, il s'ensuit aussitôt que $\symup{\underline{\operatorname{Isom}}_S(X,Y)}$
existe également, et se représente par une partie ouverte de
$\symup{\underline{\operatorname{Hom}}_S(X,Y)}$. De même, lorsque $\symup{X}$ est plat et projectif sur $\symup{S}$
et $\symup{Y}$ quasi-projectif sur $\symup{S}$, le $\symup{S}$-préschéma
$\symup{\underline{\operatorname{Imm}}_S(X/Y)}$, correspondant au sous-foncteur du foncteur
représenté par $\symup{\underline{\operatorname{Hom}}_S(X,Y)}$ qui correspond aux
$\symup{S'}$-homomorphismes $\symup{X'\to Y'}$ qui sont des immersions, est
également représentable par une partie ouverte de
$\symup{\underline{\operatorname{Hom}}_S(X,Y)}$.

Soient $\symup{\mathcal L}$ (resp. $\symup{\mathcal M}$) un faisceau inversible sur $\symup{X}$ (resp.
$\symup{Y}$) très ample relativement à $\symup{S}$, d'où un faisceau très ample
$\symup{\mathcal L\otimes_{\mathcal O_S}\mathcal M}$ sur $\symup{X\times_SY}$ relativement à
$\symup{S}$. Par suite, pour tout polynôme $\symup{P}$ à coefficients rationnels,
$\symup{\underline{\operatorname{Hilb}}^P_{(X\times_SY)/S}}$ est défini et est un préschéma
quasi-projectif sur $\symup{S}$. Il induit donc sur $\symup{\underline{\operatorname{Hom}}_S(X,Y)}$ une
partie à la fois ouverte et fermée quasi-projective sur $\symup{S}$, que nous noterons
$\symup{\underline{\operatorname{Hom}}_S(X,Y)^P}$. Donc les sections de
$\symup{\underline{\operatorname{Hom}}_S(X,Y)^P}$ sur $\symup{S}$ sont les $\symup{S}$-morphismes
$\symup{g:X\to Y}$ tels que pour tout entier $\symup{n}$, on ait
\[
\symup{\chi\!\left(\left(\mathcal L\otimes_{\mathcal O_X}g^*(\mathcal M)\right)^{
\otimes n}\right)=P(n).}
\]

On obtient de cette façon des généralisations du théorème de Matsusaka,
affirmant que les automorphismes d'une variété projective "polarisée" forment
un groupe algébrique, assertion qui prend ici une signification évidemment plus
précise, puisque nous disposons d'une définition de ce groupe comme solution
d'un problème universel. On notera d'ailleurs que, sur un corps algébriquement
clos, le groupe des automorphismes considéré anciennement est celui déduit du
"vrai" défini ici en divisant par les éléments nilpotents ; cela explique
pourquoi il y a peu de chance que les constructions plus anciennes puissent se
faire sur un corps de base non parfait, l'Idéal des éléments nilpotents
apparaissant après extension du corps de base n'étant pas nécessairement
"défini sur $\symup{k}$". Cette même remarque s'applique d'ailleurs également à la
plupart des constructions ancien style.

\src{269}
\fsection{5}{\normalfont\uline{Étude différentielle des schémas de Hilbert}.}

Elle découle du résultat suivant :
\result{PROPOSITION}{5.1} Soient $\symup{S}$ un préschéma, $\symup{S_0}$ un sous-préschéma
défini par un Idéal quasi-cohérent $\symup{\mathcal J}$ de carré nul, $\symup{X}$ un
$\symup{S}$-préschéma, $\symup{\mathcal F}$ un Module quasi-cohérent sur $\symup{X}$,
$\symup{X_0=X\times_SS_0}$ et
$\symup{\mathcal F_0=\mathcal F\otimes_{\mathcal O_S}\mathcal O_{S_0}}$, enfin
$\symup{\mathcal G_0=\mathcal F_0/\mathcal K_0}$ un Module quotient quasi-cohérent de
$\symup{\mathcal F_0}$, plat sur $\symup{S_0}$. Pour tout ouvert $\symup{U}$ de $\symup{X}$, soit
$\symup{\mathcal E(U)}$ l'ensemble des Modules quotients quasi-cohérents $\symup{\mathcal G}$ de
$\symup{\mathcal F|U}$, plats sur $\symup{S}$, et tels que
$\symup{\mathcal G\otimes_{\mathcal O_S}\mathcal O_{S_0}=\mathcal G_0}$ ; ainsi pour
$\symup{U}$ variable, les $\symup{\mathcal E(U)}$ sont les sections d'un faisceau $\symup{\mathcal E}$
sur $\symup{X}$. Ceci posé, le faisceau en groupes
\[
\symup{\mathcal A=\underline{\operatorname{Hom}}_{\mathcal O_{X_0}}
(\mathcal K_0,\mathcal G_0\otimes_{\mathcal O_{S_0}}\mathcal J)}
\]
opère de façon naturelle sur $\symup{\mathcal E}$, qui devient ainsi un faisceau
"formellement principal homogène sous $\symup{\mathcal A}$" (i. e. pour tout ouvert
$\symup{U}$ dans $\symup{X}$, $\symup{\mathcal E(U)}$ est vide ou un ensemble principal homogène sous
$\symup{\mathcal A(U)}$).

On en conclut :

\result{COROLLAIRE}{5.2} Supposons qu'il existe localement sur $\symup{X}$ un
prolongement $\symup{\mathcal G}$ de $\symup{\mathcal G_0}$ en un quotient de $\symup{\mathcal F}$ plat
sur $\symup{S}$, (i. e. que les fibres du faisceau $\symup{\mathcal E}$ soient non vides). Alors
il existe une classe d'obstruction canonique
\[
\symup{c(\mathcal G_0)\in H^1(X,\mathcal A),}
\]
dont l'annulation est nécessaire et suffisante pour l'existence d'un
prolongement global $\symup{\mathcal G}$ de $\symup{\mathcal G_0}$ en un quotient de
$\symup{\mathcal F}$ plat sur $\symup{S}$. Si cette classe est nulle, alors l'ensemble
$\symup{\mathcal E(X)}$ de tous les prolongements possibles est un ensemble principal
homogène sous $\symup{\mathcal A(X)=\operatorname{Hom}_{\mathcal O_X}(\mathcal K_0,\mathcal G_0\otimes_{\mathcal O_{S_0}}\mathcal J)}$.

L'existence du prolongement global est alors garantie en particulier si
$\symup{H^1(X,\mathcal A)=0}$.

\result{COROLLAIRE}{5.3} Supposons que
$\symup{Q=\underline{\operatorname{Quot}}_{\mathcal F/X/S}}$ existe (cf. 4. (a)) - par exemple que
$\symup{X}$ soit quasi-projectif sur $\symup{S}$ localement noethérien, et $\symup{\mathcal F}$
cohérent. Soit $\symup{x\in Q}$, correspondant à une extension résiduelle $\symup{K=k(x)}$
d'un $\symup{k(s)}$ ($\symup{s\in S}$) :
\src{270}
donc $\symup{x}$ est défini par un Module quotient cohérent
$\symup{\mathcal G_0=\mathcal F_0/\mathcal K_0}$ du Module
$\symup{\mathcal F_0=F_K}$ sur le $\symup{K}$-préschéma $\symup{X_K}$. Soit $\symup{\mathcal A}$ le
faisceau cohérent sur $\symup{X_K}$ défini par
\[
\symup{\mathcal A=\underline{\operatorname{Hom}}_{\mathcal O_{X_0}}(\mathcal K_0,\mathcal G_0).}
\]
Alors l'espace tangent de Zariski de la fibre $\symup{Q_s}$ au point $\symup{x}$ (dual sur $\symup{K}$
de $\symup{m/m^2}$, où $\symup{m}$ est l'idéal maximal de
$\symup{\mathcal O_{Q_K,x}}$) est canoniquement isomorphe à $\symup{H^0(X_K,\mathcal A)}$.

Le résultat donnant l'espace tangent de Zariski peut se généraliser, et donne
une caractérisation, pour un $\symup{S}$-morphisme donné $\symup{g:S'\to Q}$, i. e.
une section $\symup{g'}$ de $\symup{Q'=Q\times_SS'}$ sur $\symup{S'}$, du Module
\[
\symup{\Omega=g^*(\Omega^1_{Q/S})=g'^*(\mathcal J/\mathcal J^2)}
\]
(où $\symup{\mathcal J}$ est l'Idéal sur $\symup{Q'}$ défini par la section $\symup{g'}$ de $\symup{Q'}$ sur
$\symup{S'}$), par la formule fonctorielle en le Module cohérent $\symup{\mathcal M}$ sur $\symup{S'}$ :
\[
\symup{\operatorname{Hom}_{\mathcal O_{S'}}(\Omega,\mathcal M)
\simeq H^0(X',\mathcal A),}
\]
où $\symup{\mathcal A}$ est encore le Module sur $\symup{X'=X\times_SS'}$ défini par
\[
\symup{\mathcal A=\underline{\operatorname{Hom}}_{\mathcal O_{X'}}
(\mathcal K,\mathcal G\otimes_{\mathcal O_{S'}}\mathcal M),}
\]

($\symup{\mathcal G=\mathcal F'/\mathcal K}$ étant le Module quotient de
$\symup{\mathcal F'=\mathcal F\otimes_{\mathcal O_S}\mathcal O_{S'}}$ qui correspond à
$\symup{g}$). Il suffit en effet d'appliquer (5.1) en y remplaçant $\symup{S_0}$ par $\symup{S'}$, $\symup{S}$
par le préschéma $\symup{D(\mathcal M)=(S',\mathcal O_{S'}+\mathcal M)}$, où
$\symup{\mathcal M}$ est considéré comme un Idéal de carré nul.

Lorsque dans (5.1), on a $\symup{\mathcal F=\mathcal O_X}$, donc la donnée de
$\symup{\mathcal G_0}$ correspond à la donnée d'un sous-préschéma fermé $\symup{Y_0}$ de $\symup{X_0}$
plat sur $\symup{S_0}$ défini par l'Idéal $\symup{\mathcal J_0=\mathcal K_0}$, alors (5.1) donne
\[
\symup{\mathcal A=\underline{\operatorname{Hom}}_{\mathcal O_{X_0}}
(\mathcal J_0/\mathcal J_0^2,
\mathcal O_{Y_0}\otimes_{\mathcal O_{S_0}}\mathcal J),}
\]
où $\symup{\mathcal J_0/\mathcal J_0^2}$ s'interprète comme le \uline{faisceau conormal} à $\symup{Y_0}$
dans $\symup{X_0}$, qu'on note aussi
\src{271}
$\symup{\mathcal N_{Y_0/X_0}}$ ; il y a alors intérêt à considérer $\symup{\mathcal A}$
comme un Module sur $\symup{Y_0}$, et de calculer $\symup{H^0}$ et $\symup{H^1}$ sur $\symup{Y_0}$. D'ailleurs,
si $\symup{Y_0}$ est localement une intersection complète dans $\symup{X_0}$, $\symup{X}$ étant plat
sur $\symup{S}$, alors dans (5.1) la possibilité de prolongement local est garantie,
d'autre part $\symup{\mathcal J_0/\mathcal J_0^2}$ est localement libre sur $\symup{Y_0}$ et on peut
écrire
\[
\symup{\mathcal A=\check{\mathcal N}_{X_0/Y_0}\otimes_{\mathcal O_{S_0}}\mathcal J,}
\]
où le premier facteur du deuxième membre est le faisceau normal à $\symup{Y_0}$ dans
$\symup{X_0}$. Utilisant le critère fondamental de simplicité ([3], III, 3.1), on
trouve par exemple :

\result{COROLLAIRE}{5.4} Sous les conditions de (5.3), supposons que
$\symup{\mathcal F=\mathcal O_X}$, $\symup{X}$ plat sur $\symup{S}$, et que le sous-préschéma fermé
$\symup{Y_0}$ de $\symup{X_0}$ qui correspond à $\symup{\mathcal G_0}$ soit localement une intersection
complète. Alors l'espace tangent de Zariski à $\symup{Q_s}$ en $\symup{x}$ est canoniquement
isomorphe à $\symup{H^0(Y_0,\check{\mathcal N}_{X_0/Y_0})}$. Si
$\symup{H^1(Y_0,\check{\mathcal N}_{X_0/Y_0})=0}$, alors le préschéma de Hilbert
$\symup{Q}$ est simple sur $\symup{S}$ en le point $\symup{x}$.
($\symup{\check{\mathcal N}_{X_0/Y_0}}$ est le faisceau normal à $\symup{Y_0}$ dans $\symup{X_0}$.)

\result{REMARQUE}{5.5} Cet énoncé s'applique en particulier lorsque $\symup{Y_0}$ est
une intersection complète dans $\symup{X_0}$ définie par \uline{une} équation, i. e. est un
"diviseur de Cartier" positif. Alors $\symup{\check{\mathcal N}_{X_0/Y_0}}$ est isomorphe au
faisceau sur $\symup{Y_0}$ induit par le faisceau inversible $\symup{\mathcal J_0^{-1}}$ sur
$\symup{X_0}$ défini par le diviseur $\symup{Y_0}$. C'est la situation rencontrée en particulier
dans l'étude des familles de diviseurs positifs sur une variété projective non
singulière $\symup{X_0}$. L'isomorphisme entre l'espace tangent de Zariski en le point
$\symup{x}$ de $\symup{Q}$ (ou si on préfère, de l'ouvert $\symup{D}$ de $\symup{Q}$ qui correspond aux
diviseurs) et $\symup{\mathcal H^0(Y_0,\check{\mathcal N}_{X_0/Y_0})}$ était connu en géométrie
algébrique classique sous le nom de "\uline{homomorphisme caractéristique}" du premier
dans le second. Il n'était défini que lorsque $\symup{x}$ était un point simple de la
variété des paramètres $\symup{T}$ d'une "famille continue complète" de diviseurs,
c'est-à-dire, de notre point de vue, une composante irréductible du schéma
$\symup{D}$, munie de la structure \uline{réduite} induite. L'espace tangent à $\symup{T}$ en $\symup{x}$ est alors
un \uline{sous-espace} de l'espace tangent à $\symup{D}$ en $\symup{x}$, donc l'homomorphisme
caractéristique des anciens est bien injectif, mais n'est surjectif que sous des
conditions supplémentaires, par exemple si $\symup{D}$ est intègre en $\symup{x}$. En fait,
ZAPPA [8] a construit un exemple
\src{272}
(avec $\symup{X}$ surface projective non singulière sur le \uline{corps des complexes}) où même
en le point générique de $\symup{T}$, l'homomorphisme caractéristique n'est pas
surjectif. \uline{Cela signifie donc que} $\symup{D}$ \uline{n'est pas intègre même en le point
générique de la composante irréductible envisagée}. Cela montre de façon
particulièrement frappante comment les variétés à éléments nilpotents sont
nécessaires pour comprendre des phénomènes de la plus classique théorie des
surfaces.

5.6. - On a donné dans (5.4) un critère de simplicité, s'appliquant en particulier aux
schémas de diviseurs. KODAIRA a donné dans [6] un critère différent, savoir la
nullité de $\symup{H^1(X_0,\mathcal L)}$, où $\symup{\mathcal L=\mathcal J_0^{-1}}$ est le
faisceau inversible sur $\symup{X_0}$ défini par le diviseur $\symup{Y_0}$ ; critère valable
lorsque $\symup{S}$ est le spectre d'un corps de caractéristique 0, et prouvé dans [6]
par voie transcendante dans le cas où le corps de base est $\symup{\text{\uwave{C}}}$. Notons
ici que, de façon générale, $\symup{S}$ étant à nouveau quelconque, la condition de
Kodaira est une condition suffisante pour que le morphisme canonique $\symup{D\to\underline{\operatorname{Pic}}_{X/S}}$ du préschéma des diviseurs dans le préschéma de Picard de $\symup{X/S}$ soit simple au
point $\symup{x}$ envisagé (comme on le vérifie facilement par le critère habituel de
simplicité, une fois acquise l'existence de $\symup{\underline{\operatorname{Pic}}_{X/S}}$). Si
donc, de plus, $\symup{\underline{\operatorname{Pic}}_{X/S}}$ est simple sur $\symup{S}$ au point image de
$\symup{x}$ (par exemple si $\symup{\underline{\operatorname{Pic}}_{X/S}}$ est simple sur $\symup{S}$), alors $\symup{D}$
est simple sur $\symup{S}$ en $\symup{x}$. D'autre part, CARTIER a démontré que tout préschéma
en groupes localement de type fini sur \uline{un corps} $\symup{k}$ \uline{de caractéristique} 0 est
simple sur $\symup{k}$. En conjuguant ces deux résultats, on retrouve le résultat de
KODAIRA. On notera qu'il résulte de ces remarques que sur un corps $\symup{K}$ de
caractéristique $\symup{p>0}$, lorsque $\symup{\underline{\operatorname{Pic}}_{X/S}}$ n'est pas simple sur
$\symup{K}$ (ce qui est le cas lorsque $\symup{X}$ est la surface de Igusa), la condition
$\symup{H^1(X_0,\mathcal L)=0}$ implique au contraire que $\symup{D}$ n'est pas simple en $\symup{x}$,
et même n'est pas réduit en $\symup{x}$ si $\symup{K}$ est algébriquement clos.

Pour finir, signalons encore le résultat suivant, qui joue un rôle important
dans l'étude différentielle des espaces fibrés :

\result{PROPOSITION}{5.7} Soient $\symup{X}$ un préschéma fini et plat sur $\symup{S}$
localement noethérien, et soit $\symup{Z}$ un préschéma sur $\symup{X}$, tel que
$\symup{\prod\limits_{X/S}(Z/X)}$ existe (ce qui est le cas si $\symup{Z}$ est quasi-projectif sur $\symup{X}$).
Si alors $\symup{Z}$ est simple sur $\symup{X}$, $\symup{\prod\limits_{X/S}(Z/X)}$ est simple sur $\symup{S}$.

C'est une conséquence immédiate de la définition, et du critère habituel de
simplicité ([3], III, 3.1). Remarquons que lorsque $\symup{X}$ est fini et plat sur
$\symup{S}$, la question de l'existence de $\symup{\prod\limits_{X/S}(Z/X)}$ peut se traiter très élémentairement, sans
\src{273}
utiliser la théorie des schémas de Hilbert. On trouve par exemple que si $\symup{X}$
est radiciel sur $\symup{S}$, alors $\symup{\prod\limits_{X/S}(Z/X)}$ existe sans aucune restriction sur
$\symup{Z}$. Par exemple, soit $\symup{T}$ un $\symup{S}$-préschéma, et soit $\symup{T_n}$ le "voisinage
infinitésimal d'ordre $\symup{n}$" de la diagonale de $\symup{T\times_ST}$ dans $\symup{T\times_ST}$,
muni des morphismes $\symup{p_1,p_2:T_n\to T}$ induits par les deux
projections. On considère $\symup{T_n}$ comme un préschéma fini sur $\symup{T}$ grâce à $\symup{p_1}$,
supposons-le de plus plat sur $\symup{T}$ (ce qui est le cas si $\symup{T}$ est simple sur $\symup{S}$).
Pour tout préschéma $\symup{X}$ sur $\symup{T}$, posons
\[
\symup{(X/T/S)^{(n)}=
\prod\limits_{T_n/T}\bigl(p_2^*(X/T)/T_n\bigr),}
\]

c'est un préschéma sur $\symup{T}$ appelé \uline{fibré des germes de sections d'ordre}
$\symup{n}$ \uline{de} $\symup{X}$ \uline{sur} $\symup{T}$ (relativement à $\symup{S}$). Il dépend fonctoriellement de $\symup{X}$, et
est simple sur $\symup{T}$ si $\symup{X}$ l'est.

\fsection{6}{\normalfont\uline{Relations avec la notion de norme et les produits symétriques}.}

Soient $\symup{S}$ un préschéma, $\symup{X}$ et $\symup{Y}$ des $\symup{S}$-préschémas,
\[
\symup{u:(X/S)^n\to Y}
\]
un $\symup{S}$-morphisme \uline{symétrique} de la puissance cartésienne $\symup{n}$-ième de $\symup{X/S}$ dans
$\symup{Y}$. Nous supposerons pour simplifier $\symup{S}$ localement noethérien, et $\symup{X}$ et $\symup{Y}$
de type fini sur $\symup{S}$. On peut alors, à tout Module cohérent $\symup{\mathcal F}$ sur
$\symup{X}$, à support fini sur $\symup{S}$, plat sur $\symup{S}$ et de rang relatif sur $\symup{S}$ égal à $\symup{n}$
(i. e. tel que $\symup{f_*(\mathcal F)}$ soit un Module localement libre de rang $\symup{n}$ sur
$\symup{S}$), faire correspondre de façon naturelle une section de $\symup{Y}$ sur $\symup{S}$ :
\[
\symup{\mathcal N^u_{X/S}(\mathcal F)\in\Gamma(Y/S).}
\]
Nous ne donnerons pas ici la définition en forme, nous contentant de signaler
que le formalisme auquel on arrive est une généralisation naturelle du
formalisme habituel des normes et traces. Lorsque la puissance symétrique $\symup{n}$-ième de $\symup{X}$ sur $\symup{S}$ existe (par exemple lorsque les orbites du groupe symétrique $\symup{\mathfrak S_n}$ opérant sur $\symup{(X/S)^n}$ sont contenues dans des ouverts affines), on peut prendre pour $\symup{Y}$ cette puissance symétrique $\symup{\operatorname{Symm}^n_S(X)}$, et on trouve un élément canonique
\[
\symup{\mathcal N_{X/S}(\mathcal F)\in\Gamma(\operatorname{Symm}^n_S(X)/S),}
\]
\src{274}
qui permet de retrouver les $\symup{\mathcal N^u_{X/S}(F)}$. Un autre cas important
est celui où $\symup{X}$ est un monoïde commutatif sur $\symup{S}$, et $\symup{X=Y}$, le morphisme $\symup{u}$
provenant de la loi de composition de $\symup{X}$. On écrira alors simplement
$\symup{\mathcal N(\mathcal F)}$ pour la section de $\symup{X}$ sur $\symup{S}$ associée au Module
$\symup{\mathcal F}$ sur $\symup{X}$.

Supposons maintenant que l'on ait un Module cohérent $\symup{\mathcal F}$ sur $\symup{X}$, tel
que $\symup{\underline{\operatorname{Quot}}_{\mathcal F/X/S}}$ existe, ou du moins tel que le
foncteur $\symup{\mathop{\text{\uwave{Quot}}}\nolimits^n_{\mathcal F/X/S}}$ qui associe à
tout $\symup{S'}$ sur $\symup{S}$ l'ensemble des faisceaux quotients cohérents $\symup{\mathcal M}$ de
$\symup{\mathcal F'=\mathcal F\otimes_{\mathcal O_S}\mathcal O_{S'}}$ qui sont plats
sur $\symup{S'}$ et de rang relatif $\symup{n}$, soit représentable par un $\symup{S}$-préschéma
$\symup{\underline{\operatorname{Quot}}^n_{\mathcal F/X/S}}$ (Lorsque $\symup{X}$ est quasi-projectif sur
$\symup{S}$, alors $\symup{\underline{\operatorname{Quot}}^n_{\mathcal F/X/S}}$ existe bien, et n'est
autre, avec les notations du numéro 3, que
$\symup{\underline{\operatorname{Quot}}^P_{\mathcal F/X/S}}$, où $\symup{P}$ est le polynôme réduit au
terme constant $\symup{n}$.) Comme la formation des $\symup{\mathcal N^u_{X/S}(\mathcal M)}$ est
compatible avec le changement de base, on obtient alors un morphisme canonique
\[
\symup{\mathcal N^u_{X/S}:\underline{\operatorname{Quot}}^n_{\mathcal F/X/S}\to Y,}
\]
en particulier, si la puissance symétrique $\symup{n}$-ième de $\symup{X}$ sur $\symup{S}$ existe :
\[
\symup{\mathcal N_{X/S}:\underline{\operatorname{Quot}}^n_{\mathcal F/X/S}\to\operatorname{Symm}^n_S(X).}
\]
Le cas le plus important est celui où $\symup{\mathcal F=\mathcal O_X}$, qui donne un morphisme :
\[
\symup{\mathcal N_{X/S}:\underline{\operatorname{Hilb}}^n_{X/S}
\to\operatorname{Symm}^n_S(X).}
\]
C'est évidemment un isomorphisme pour $\symup{n=0}$ ou $\symup{n=1}$. Mais pour
$\symup{n>1}$, même
si $\symup{S}$ est le spectre d'un corps $\symup{k}$, et lorsque $\symup{X}$ est simple sur $\symup{S}$, ce
n'est pas en général un isomorphisme ni même un morphisme injectif, puisque un
sous-schéma de dimension 0 de $\symup{X}$ (correspondant par exemple à un idéal
$\symup{i}$ primaire pour l'idéal maximal dans un anneau local
$\symup{\mathcal O_{X,x}}$, pour un point fermé $\symup{x}$ de $\symup{X}$) n'est pas connu quand on
connaît le cycle qu'il définit (en l'occurrence, quand on connaît la codimension
sur $\symup{k}$ de $\symup{i}$ dans $\symup{\mathcal O_{X,x}}$). On peut seulement dire ceci
($\symup{S}$ étant de nouveau quelconque) :

\medskip
\noindent a. Si $\symup{X}$ est simple sur $\symup{S}$, alors le morphisme-norme définit un
isomorphisme de l'ouvert de $\symup{\underline{\operatorname{Hilb}}^n_{X/S}}$ qui correspond à la
classification de revêtements étales de rang $\symup{n}$ contenus dans $\symup{X}$ (cf.
n\textsuperscript{o} 4, (b)), avec l'ouvert de
$\symup{\operatorname{Symm}^n_S(X)}$ qui correspond aux $\symup{n}$-cycles sans composantes
multiples.
\src{275}

\medskip
\noindent b. Si de plus $\symup{X}$ est de dimension relative 1 sur $\symup{S}$, alors le
morphisme norme définit même un isomorphisme de
$\symup{\underline{\operatorname{Hilb}}^n_{X/S}}$ avec $\symup{\operatorname{Symm}^n_{X/S}}$.

Ce dernier fait est dû à ce qu'un sous-schéma de dimension 0 d'une courbe
algébrique non singulière est connu quand on connaît le cycle correspondant. La
même remarque s'applique d'ailleurs plus généralement aux diviseurs de Cartier,
positifs sur un schéma algébrique non singulier (et il n'est pas exclu que dans
ce cas très particulier, la variété de Chow donne la même chose que le schéma
de Hilbert).

\fsection{7}{\normalfont\uline{Compléments et questions}.}

Comme l'a remarqué J.-P. SERRE, il résulte d'un exemple bien connu de NAGATA
qu'on peut trouver un schéma $\symup{S}$, spectre d'un corps $\symup{k}$, un $\symup{S}$-schéma $\symup{S'}$,
spectre d'une extension quadratique $\symup{k'}$ de $\symup{k}$, enfin un $\symup{S'}$-schéma $\symup{X}$
propre et simple (mais non projectif) de dimension 3, tel que
$\symup{\prod\limits_{S'/S}(X/S')}$ n'existe pas. Cela implique a fortiori que le schéma de
Hilbert $\symup{\underline{\operatorname{Hilb}}^2_{X/S}}$ n'existe pas (ni même le $\symup{k}$-schéma qui
représenterait les revêtements étales de rang 2 de $\symup{S}$ contenus dans $\symup{X}$, ni a
fortiori le carré symétrique de $\symup{X}$, cf. numéro précédent). Cela impose donc
des limitations sérieuses aux possibilités de constructions non projectives en
Géométrie algébrique. (Il est cependant plausible que de telles limitations ne
se présenteront pas en Géométrie analytique, pas plus qu'elles ne se présentent
en Géométrie formelle, (cf. [2], II)). Par contre, si $\symup{X}$ est un schéma propre
sur le spectre $\symup{S}$ d'un corps $\symup{k}$, et si $\symup{Z}$ est quasi-projectif sur $\symup{X}$, alors
$\symup{\prod\limits_{X/S}(Z/X)}$ existe, et est un schéma, somme d'une suite de schémas
quasi-projectifs sur $\symup{S}$ (comme dans le cas projectif (3.1)). Pour le voir, on
se ramène en effet au cas où $\symup{X}$ est lui-même projectif, en dominant $\symup{X}$ par un
$\symup{S}$-schéma projectif $\symup{X'}$ ; nous ne donnerons pas ici le détail de la
démonstration, qui utilise aussi le résultat de factorisation d'un morphisme
fini signalé dans [2], I, A, n\textsuperscript{o} 2 (b). Le succès de la méthode
tient au fait que, $\symup{S}$ étant le spectre d'un corps, le $\symup{X'}$ qui intervient dans
le lemme de Chow sera automatiquement plat sur $\symup{S}$.
J'ignore si le résultat reste valable sans hypothèse sur $\symup{S}$, en supposant
seulement $\symup{X}$ propre et plat sur $\symup{S}$, $\symup{Z}$ quasi-projectif sur $\symup{X}$. Un cas
important dans les applications est celui où $\symup{Z}$ est un sous-schéma fermé de
$\symup{X}$ ; si alors $\symup{\prod\limits_{X/S}(Z/X)}$ existe, c'est nécessairement un sous-schéma
fermé de $\symup{S}$. On peut le
\src{276}
construire directement de façon assez simple lorsque $\symup{X}$ est projectif sur $\symup{S}$,
sans utiliser la théorie des schémas de Hilbert, et la méthode employée montre
plus généralement que si $\symup{Z}$ est affine sur $\symup{X}$, alors $\symup{\prod\limits_{X/S}(Z/X)}$ existe
et est affine sur $\symup{S}$. Elle montre également que si $\symup{X}$ est propre et plat sur
$\symup{S}$ (pas nécessairement projectif sur $\symup{S}$), alors pour tout fibré vectoriel $\symup{Z}$
localement trivial sur $\symup{X}$, $\symup{\prod\limits_{X/S}(Z/X)}$ existe, et est un fibré vectoriel
sur $\symup{S}$. Il serait désirable que ces résultats soient repris et unifiés.

\begin{center}BIBLIOGRAPHIE\end{center}
\addcontentsline{toc}{section}{Bibliographie}

\begin{description}
\item[{\normalfont[1]}] BOREL (A.) et SERRE (J.-P.). - Le théorème de Riemann-Roch,
Bull. Soc. math. France, t. 86, 1958, p. 97-136.

\item[{\normalfont[2]}] GROTHENDIECK (Alexander). - Technique de construction et
théorèmes d'existence en géométrie algébrique, I., II., III., Séminaire
Bourbaki, t. 13, 1959/60, n\textsuperscript{o} 190, 29 p.,
n\textsuperscript{o} 195, 22 p. ; et t. 14, 1960/61,
n\textsuperscript{o} 212, 20 p.

\item[{\normalfont[3]}] GROTHENDIECK (Alexander). - Séminaire de Géométrie algébrique,
I., II., III., IV. - Paris, Institut des hautes Études scientifiques, 1960/61
(multigraphié).

\item[{\normalfont[4]}] GROTHENDIECK (Alexander). - Technique de construction en
géométrie analytique, IV., V., Séminaire Cartan, t. 13, 1960/61,
n\textsuperscript{o} 11 et 12.

\item[{\normalfont[5]}] GROTHENDIECK (A.) et DIEUDONNÉ (J.). - Éléments de
géométrie algébrique, I : Le langage des schémas. - Paris, Presses
universitaires de France, 1960 (Institut des hautes Études scientifiques,
Publications mathématiques, 4) ; II., III., IV. (à paraître).

\item[{\normalfont[6]}] KODAIRA (K.). - Characteristic linear systems of complete
continuous systems, Amer. J. of Math., t. 78, 1956, p. 716-744.

\item[{\normalfont[7]}] SERRE (Jean-Pierre). - Exemples de variétés projectives en
caractéristique $\symup{p}$ non relevables en caractéristique 0, Proc. Nat. Acad.
Sc. U. S. A., t. 47, 1961, p. 108-109.

\item[{\normalfont[8]}] ZAPPA (G.). - Sull'esistenza, sopra le superficie algebriche, di
sistemi continui completi infiniti, la cui curva è a serie caratteristica
incompleta, Pont. Acad. Sc. Acta, t. 9, 1945, p. 91-93.
\end{description}

\begin{center}ADDITIF\end{center}
\addcontentsline{toc}{section}{Additif}
\begin{center}
[ajouté à la correction des épreuves]
\end{center}

Il apparaît maintenant que les conjectures de [2], III,
n\textsuperscript{o} 8, sont fausses, même pour des variétés non singulières sur
un corps de caractéristique 0, tant en ce qui concerne l'existence que la
quasi-projectivité du quotient, et même lorsque $\symup{\mathfrak G}$ opère avec un graphe fermé.

\begin{center}
  \rule{.14\linewidth}{.3pt}
\end{center}
